% Options for packages loaded elsewhere
\PassOptionsToPackage{unicode}{hyperref}
\PassOptionsToPackage{hyphens}{url}
\PassOptionsToPackage{dvipsnames,svgnames,x11names}{xcolor}
%
\documentclass[
  10pt,
]{article}
\usepackage{amsmath,amssymb}
\usepackage{iftex}
\ifPDFTeX
  \usepackage[T1]{fontenc}
  \usepackage[utf8]{inputenc}
  \usepackage{textcomp} % provide euro and other symbols
\else % if luatex or xetex
  \usepackage{unicode-math} % this also loads fontspec
  \defaultfontfeatures{Scale=MatchLowercase}
  \defaultfontfeatures[\rmfamily]{Ligatures=TeX,Scale=1}
\fi
\usepackage{lmodern}
\ifPDFTeX\else
  % xetex/luatex font selection
\fi
% Use upquote if available, for straight quotes in verbatim environments
\IfFileExists{upquote.sty}{\usepackage{upquote}}{}
\IfFileExists{microtype.sty}{% use microtype if available
  \usepackage[]{microtype}
  \UseMicrotypeSet[protrusion]{basicmath} % disable protrusion for tt fonts
}{}
\makeatletter
\@ifundefined{KOMAClassName}{% if non-KOMA class
  \IfFileExists{parskip.sty}{%
    \usepackage{parskip}
  }{% else
    \setlength{\parindent}{0pt}
    \setlength{\parskip}{6pt plus 2pt minus 1pt}}
}{% if KOMA class
  \KOMAoptions{parskip=half}}
\makeatother
\usepackage{xcolor}
\usepackage{color}
\usepackage{fancyvrb}
\newcommand{\VerbBar}{|}
\newcommand{\VERB}{\Verb[commandchars=\\\{\}]}
\DefineVerbatimEnvironment{Highlighting}{Verbatim}{commandchars=\\\{\}}
% Add ',fontsize=\small' for more characters per line
\newenvironment{Shaded}{}{}
\newcommand{\AlertTok}[1]{\textcolor[rgb]{1.00,0.00,0.00}{\textbf{#1}}}
\newcommand{\AnnotationTok}[1]{\textcolor[rgb]{0.38,0.63,0.69}{\textbf{\textit{#1}}}}
\newcommand{\AttributeTok}[1]{\textcolor[rgb]{0.49,0.56,0.16}{#1}}
\newcommand{\BaseNTok}[1]{\textcolor[rgb]{0.25,0.63,0.44}{#1}}
\newcommand{\BuiltInTok}[1]{\textcolor[rgb]{0.00,0.50,0.00}{#1}}
\newcommand{\CharTok}[1]{\textcolor[rgb]{0.25,0.44,0.63}{#1}}
\newcommand{\CommentTok}[1]{\textcolor[rgb]{0.38,0.63,0.69}{\textit{#1}}}
\newcommand{\CommentVarTok}[1]{\textcolor[rgb]{0.38,0.63,0.69}{\textbf{\textit{#1}}}}
\newcommand{\ConstantTok}[1]{\textcolor[rgb]{0.53,0.00,0.00}{#1}}
\newcommand{\ControlFlowTok}[1]{\textcolor[rgb]{0.00,0.44,0.13}{\textbf{#1}}}
\newcommand{\DataTypeTok}[1]{\textcolor[rgb]{0.56,0.13,0.00}{#1}}
\newcommand{\DecValTok}[1]{\textcolor[rgb]{0.25,0.63,0.44}{#1}}
\newcommand{\DocumentationTok}[1]{\textcolor[rgb]{0.73,0.13,0.13}{\textit{#1}}}
\newcommand{\ErrorTok}[1]{\textcolor[rgb]{1.00,0.00,0.00}{\textbf{#1}}}
\newcommand{\ExtensionTok}[1]{#1}
\newcommand{\FloatTok}[1]{\textcolor[rgb]{0.25,0.63,0.44}{#1}}
\newcommand{\FunctionTok}[1]{\textcolor[rgb]{0.02,0.16,0.49}{#1}}
\newcommand{\ImportTok}[1]{\textcolor[rgb]{0.00,0.50,0.00}{\textbf{#1}}}
\newcommand{\InformationTok}[1]{\textcolor[rgb]{0.38,0.63,0.69}{\textbf{\textit{#1}}}}
\newcommand{\KeywordTok}[1]{\textcolor[rgb]{0.00,0.44,0.13}{\textbf{#1}}}
\newcommand{\NormalTok}[1]{#1}
\newcommand{\OperatorTok}[1]{\textcolor[rgb]{0.40,0.40,0.40}{#1}}
\newcommand{\OtherTok}[1]{\textcolor[rgb]{0.00,0.44,0.13}{#1}}
\newcommand{\PreprocessorTok}[1]{\textcolor[rgb]{0.74,0.48,0.00}{#1}}
\newcommand{\RegionMarkerTok}[1]{#1}
\newcommand{\SpecialCharTok}[1]{\textcolor[rgb]{0.25,0.44,0.63}{#1}}
\newcommand{\SpecialStringTok}[1]{\textcolor[rgb]{0.73,0.40,0.53}{#1}}
\newcommand{\StringTok}[1]{\textcolor[rgb]{0.25,0.44,0.63}{#1}}
\newcommand{\VariableTok}[1]{\textcolor[rgb]{0.10,0.09,0.49}{#1}}
\newcommand{\VerbatimStringTok}[1]{\textcolor[rgb]{0.25,0.44,0.63}{#1}}
\newcommand{\WarningTok}[1]{\textcolor[rgb]{0.38,0.63,0.69}{\textbf{\textit{#1}}}}
\usepackage{longtable,booktabs,array}
\usepackage{calc} % for calculating minipage widths
% Correct order of tables after \paragraph or \subparagraph
\usepackage{etoolbox}
\makeatletter
\patchcmd\longtable{\par}{\if@noskipsec\mbox{}\fi\par}{}{}
\makeatother
% Allow footnotes in longtable head/foot
\IfFileExists{footnotehyper.sty}{\usepackage{footnotehyper}}{\usepackage{footnote}}
\makesavenoteenv{longtable}
\setlength{\emergencystretch}{3em} % prevent overfull lines
\providecommand{\tightlist}{%
  \setlength{\itemsep}{0pt}\setlength{\parskip}{0pt}}
\setcounter{secnumdepth}{5}

\usepackage{amsmath,amssymb}
\usepackage{graphicx}
\usepackage{longtable,booktabs,array}
\usepackage{xcolor}
\usepackage[margin=2.0cm]{geometry}
\usepackage{microtype}
\setlength{\emergencystretch}{3em}
\setcounter{MaxMatrixCols}{20}
\pdfinfoomitdate=1
\pdftrailerid{}
\pdfsuppressptexinfo=-1
\renewcommand{\arraystretch}{1.12}
\usepackage{fvextra}
\RecustomVerbatimEnvironment{verbatim}{Verbatim}{breaklines,breakanywhere,fontsize=\small}
\fvset{breaklines,breakanywhere}
\allowdisplaybreaks
\sloppy
\ifLuaTeX
  \usepackage{selnolig}  % disable illegal ligatures
\fi
\IfFileExists{bookmark.sty}{\usepackage{bookmark}}{\usepackage{hyperref}}
\IfFileExists{xurl.sty}{\usepackage{xurl}}{} % add URL line breaks if available
\urlstyle{same}
\hypersetup{
  pdftitle={UM-4 provenance},
  pdfauthor={Underlying theory: Patrick L. Nash; Recorded by Claude (Anthropic); Provenance of UM-4, 2026-10-25; Repository once-ere/Dirac\_claude, folder universe\_model\_4},
  colorlinks=true,
  linkcolor={blue},
  filecolor={Maroon},
  citecolor={Blue},
  urlcolor={blue},
  pdfcreator={LaTeX via pandoc}}

\title{UM-4 provenance}
\usepackage{etoolbox}
\makeatletter
\providecommand{\subtitle}[1]{% add subtitle to \maketitle
  \apptocmd{\@title}{\par {\large #1 \par}}{}{}
}
\makeatother
\subtitle{Where every input of UM-4 comes from, every decision taken,
every computation and its checks, the external data and software, and
how to reproduce all of it}
\author{Underlying theory: Patrick L. Nash \and Recorded by Claude
(Anthropic) \and Provenance of UM-4, 2026-10-25 \and Repository
once-ere/Dirac\_claude, folder universe\_model\_4}
\date{}

\begin{document}
\maketitle

{
\hypersetup{linkcolor=}
\setcounter{tocdepth}{3}
\tableofcontents
}
\hypertarget{the-fold-definition-and-examples}{%
\section{The fold: definition and
examples}\label{the-fold-definition-and-examples}}

\textbf{Definition.} The \emph{fold} is the seven-dimensional surface of
the author's eight-dimensional spacetime on which the cosine of the
angle \(z = 6Hx_0\) of the hidden direction vanishes:

\[
\mathcal{F} = \big\{X = (x_0, x_1, \dots, x_7):\ \cos(6Hx_0) = 0\big\},\qquad
\text{in one period } 0 < z < \pi\text{ of the metric: } z = \frac{\pi}{2},\quad x_0 = \frac{\pi}{12H} \tag{FD-1}
\] \textbf{Symbols in (FD-1).}

\begin{itemize}
\tightlist
\item
  \(\mathcal{F}\): The fold: the seven-dimensional surface
  \(\cos(6Hx_0) = 0\) of the eight-dimensional spacetime (FD-1), in one
  period \(0 < z < \pi\) the surface \(z = \pi/2\), \(x_0 = \pi/(12H)\),
  where \(\sigma = 0\); \(\vert_{\mathcal{F}}\) means restricted to the
  fold. Status: DERIVED.
\item
  \(X = \{x_0, \dots, x_7\}\): The author's coordinates of the
  eight-dimensional spacetime, \(X = \{x_0, x_1, \dots, x_7\}\) with the
  hidden direction first (Dirac\_Einstein-Lovelock\_4+4\_Nash.nb,
  section DEFCHART); \(X^M\) is the coordinate with the index \(M\) and
  \(\mathrm{d}^8X = \mathrm{d}x_0\,\mathrm{d}x_1\cdots\mathrm{d}x_7\)
  the coordinate volume element. Unit: length (each coordinate). Status:
  AUTHOR.
\item
  \(x_0\): The hidden space-like coordinate (listed first); the metric
  depends on it through \(z = 6Hx_0\), periodically with the period
  \(\pi/(6H)\). Unit: length. Status: AUTHOR.
\item
  \(x_1\): The first coordinate of ordinary 3-space. Unit: length.
  Status: AUTHOR.
\item
  \(x_7\): The third of the three extra time-like coordinates. Unit:
  length. Status: AUTHOR.
\item
  \(H\): The author's constant (an inverse length) that fixes the period
  \(\pi/(6H)\) of \(x_0\), the warp factor
  \(\mathrm{e}^{-2H\lvert\sigma\rvert}\) of the metric, the cosmological
  constant \(\Lambda = -21H^2\) and the tension
  \(\kappa\lambda_{\mathrm{F}} = 12H\) of the fold; in UM-4 a free
  parameter, bounded from below by the propagation of the gauge fields
  in the bulk. Unit: GeV (1/length). Value: \(H \gtrsim 2.5\) TeV (from
  the invisible width of the Z boson). Status: AUTHOR.
\item
  \(z\): The angle variable of the hidden direction, \(z = 6Hx_0\); one
  period of \(x_0\) is \(0 < z < \pi\), the fold is \(z = \pi/2\). Unit:
  dimensionless. Status: AUTHOR.
\item
  \(\pi\): The number pi, the ratio of the circumference of a circle to
  its diameter. Unit: dimensionless. Value: 3.141592654. Status:
  STANDARD.
\item
  \(\Lambda\): The eight-dimensional cosmological constant of UM-4,
  \(\Lambda = -21H^2\): the value for which the static block is a vacuum
  of Einstein gravity (its Einstein tensor is \(21H^2\) in every
  direction). The author writes the pure number \(\Lambda_{\mathrm{N}}\)
  in
  \(0 = -\Lambda_{\mathrm{N}} + \sum_p H^{-2p}w_{(p)}\mathcal{E}_{(p)}\),
  so that \(\Lambda = -\Lambda_{\mathrm{N}}H^2\). Unit: GeV\(^2\).
  Value: \(-21H^2\); the author's original \(-10H^2\); UM-3: \(-7H^2\).
  Status: AUTHOR (the term), DERIVED (the value).
\item
  \(\Lambda_{\mathrm{N}}\): The author's pure number in the vacuum
  equations of his notebook, \(\Lambda = -\Lambda_{\mathrm{N}}H^2\): his
  original value 10; 21 in UM-4 (Einstein gravity), 7 in UM-3. Unit:
  dimensionless. Value: 21 (the author's original 10). Status: AUTHOR.
\item
  \(w_{(p)}\): The Lovelock weights of the author's Einstein-Lovelock
  action, \(w_{(1)} = 1\) and \(w_{(2)}\), \(w_{(3)}\) of the
  Gauss-Bonnet and the cubic densities. UM-4 sets
  \(w_{(2)} = w_{(3)} = 0\) (Einstein gravity): the only weights for
  which gravitons propagate on the light cone on every background, not
  only around the static block. The author's original weights
  \(w_{(2)} = 1/32\), \(w_{(3)} = 1/1152\)
  (Dirac\_Einstein-Lovelock\_4+4\_Nash.nb, In{[}134{]}) keep the
  gravitons on the light cone only around the static block (with
  \(\Lambda = -161H^2/16\)); UM-3 used \((1, 1/20, 1/360)\). Unit:
  dimensionless. Value: \((1, 0, 0)\) in UM-4; the author's original
  \((1, 1/32, 1/1152)\). Status: AUTHOR (the form), DERIVED (the values
  that the order of 2026-10-25 selects).
\end{itemize}

\(X = \{x_0, \dots, x_7\}\) are the author's coordinates: \(x_0\) the
hidden space-like direction, \(x_1, x_2, x_3\) ordinary space, \(x_4\)
the time, \(x_5, x_6, x_7\) the three extra times; \(H\) is the author's
constant (an inverse length). The metric of UM-4 depends on \(x_0\) only
through \(\cot^2z\) and \((\sin^2z)^{1/6}\), so it repeats with the
period \(\pi/(6H)\) of \(x_0\), and every period \(0 < z < \pi\)
contains exactly one surface of \(\mathcal{F}\), at \(z = \pi/2\):
\(\mathcal{F}\) is the union of these copies, one per period (example
2), and ``the fold'' without a qualifier is the copy in the period that
contains our world. A point of the fold is labelled by the seven
coordinates \(x_1, \dots, x_7\), and the metric restricted to the fold
(\(\mathrm{d}x_0 = 0\), \(z = \pi/2\)) is

\[
g_{MN}\,\mathrm{d}X^M\mathrm{d}X^N\big\vert_{\mathcal{F}} = \mathrm{e}^{2a_4}\big(\mathrm{d}x_1^2 + \mathrm{d}x_2^2 + \mathrm{d}x_3^2\big)
 - \mathrm{d}x_4^2 - \mathrm{e}^{-2a_4}b_4^2\big(\mathrm{d}x_5^2 + \mathrm{d}x_6^2 + \mathrm{d}x_7^2\big) \tag{FD-2}
\] \textbf{Symbols in (FD-2).}

\begin{itemize}
\tightlist
\item
  \(g_{MN}\): The metric tensor of the eight-dimensional spacetime
  (signature \(+,+,+,+,-,-,-,-\) in the order \(x_0, \dots, x_7\)):
  \(g_{MN}\,\mathrm{d}X^M\mathrm{d}X^N\) is the squared interval,
  \(g^{MN}\) its inverse, \(\lvert g\rvert\) the absolute value of its
  determinant (\(\sqrt{\lvert g\rvert}\,\mathrm{d}^8X\) is the proper
  volume element). Unit: dimensionless. Status: AUTHOR.
\item
  \(M\): A coordinate index of eight dimensions (also with a numbered
  copy \(M_1, M_2, \dots\)): it runs over \(x_0, x_1, \dots, x_7\); a
  repeated index is summed. Status: STANDARD.
\item
  \(N\): A coordinate index of eight dimensions (also
  \(N_1, N_2, \dots\)), running over \(x_0, \dots, x_7\); a repeated
  index is summed. Status: STANDARD.
\item
  \(\mathrm{d}\): The differential: \(\mathrm{d}\) in front of a
  quantity is an infinitesimal change of it (as \(\mathrm{d}\sigma\), a
  small step of \(\sigma\)); the quotient of two such changes is a
  derivative; \(\mathrm{d}\) with an exponent, in front of several
  variables, is their volume element, the product of their
  differentials. Status: STANDARD.
\item
  \(X = \{x_0, \dots, x_7\}\): The author's coordinates of the
  eight-dimensional spacetime, \(X = \{x_0, x_1, \dots, x_7\}\) with the
  hidden direction first (Dirac\_Einstein-Lovelock\_4+4\_Nash.nb,
  section DEFCHART); \(X^M\) is the coordinate with the index \(M\) and
  \(\mathrm{d}^8X = \mathrm{d}x_0\,\mathrm{d}x_1\cdots\mathrm{d}x_7\)
  the coordinate volume element. Unit: length (each coordinate). Status:
  AUTHOR.
\item
  \(\mathcal{F}\): The fold: the seven-dimensional surface
  \(\cos(6Hx_0) = 0\) of the eight-dimensional spacetime (FD-1), in one
  period \(0 < z < \pi\) the surface \(z = \pi/2\), \(x_0 = \pi/(12H)\),
  where \(\sigma = 0\); \(\vert_{\mathcal{F}}\) means restricted to the
  fold. Status: DERIVED.
\item
  \(\mathrm{e}\): Euler's number, the base of the natural logarithm;
  \(\mathrm{e}^{X}\) is the exponential of \(X\). Unit: dimensionless.
  Value: 2.718281828. Status: STANDARD.
\item
  \(a_4\): The author's metric function of the time \(x_4\) that scales
  ordinary 3-space: on the fold \(\mathrm{e}^{a_4}\) is the scale factor
  of cosmology and \(a_4'\) the Hubble rate. Unit: dimensionless.
  Status: AUTHOR.
\item
  \(b_4\): The author's metric function of the time \(x_4\) in the
  extra-time components, \(b_4 > 0\), never identically 1 (the author's
  order of 2026-10-24); the scale factor (the factor of the lengths) of
  the extra times is \(\mathrm{e}^{-a_4}b_4\) on the fold and
  \(\mathrm{e}^{-a_4}b_4\,\mathrm{e}^{-H\lvert\sigma\rvert}\) at the
  distance \(\sigma\) from it (the square root of the warp factor times
  \(\mathrm{e}^{-a_4}b_4\)). Unit: dimensionless. Status: AUTHOR.
\item
  \(x_1\): The first coordinate of ordinary 3-space. Unit: length.
  Status: AUTHOR.
\item
  \(x_2\): The second coordinate of ordinary 3-space. Unit: length.
  Status: AUTHOR.
\item
  \(x_3\): The third coordinate of ordinary 3-space. Unit: length.
  Status: AUTHOR.
\item
  \(x_4\): The time; a prime on a function of \(x_4\) alone is
  \(\mathrm{d}/\mathrm{d}x_4\). Unit: length (= time, \(c = 1\)).
  Status: AUTHOR.
\item
  \(x_5\): The first of the three extra time-like coordinates. Unit:
  length. Status: AUTHOR.
\item
  \(x_6\): The second of the three extra time-like coordinates. Unit:
  length. Status: AUTHOR.
\item
  \(x_7\): The third of the three extra time-like coordinates. Unit:
  length. Status: AUTHOR.
\item
  \({}'\): The derivative of a function of the time \(x_4\) alone with
  respect to \(x_4\) (\(a_4' = \mathrm{d}a_4/\mathrm{d}x_4\); two primes
  the second derivative); on a function of \(\sigma\) alone (a profile)
  the derivative with respect to \(\sigma\). Status: STANDARD.
\end{itemize}

with three space-like and four time-like directions; \(a_4\) and \(b_4\)
are the author's two functions of the time \(x_4\). Notation: \(\ln\) is
the natural logarithm, \(\mathrm{e}\) Euler's number,
\(\exp(y) = \mathrm{e}^y\), \(\sin\), \(\cos\) and \(\cot = \cos/\sin\)
the trigonometric functions, \(\sinh\) and \(\tanh\) the hyperbolic
functions, \(\lvert y\rvert\) the absolute value; \(f(0^+)\) and
\(f(0^-)\) are the limits of \(f(\sigma)\) as \(\sigma\) approaches 0
from above and from below.

\textbf{Why this surface is called a fold.} Measured by the proper
distance along the hidden direction,

\[
\sigma = \mathrm{sgn}\Big(\frac{\pi}{2} - z\Big)\frac{\ln\sin z}{6H},\qquad \mathrm{d}\sigma = \lvert\cot z\rvert\,\mathrm{d}x_0,\qquad
\sin z = \mathrm{e}^{-6H\lvert\sigma\rvert} \tag{FD-3}
\] \textbf{Symbols in (FD-3).}

\begin{itemize}
\tightlist
\item
  \(\sigma\): The proper coordinate of the hidden direction:
  \(\mathrm{d}\sigma = \lvert\cot z\rvert\,\mathrm{d}x_0\),
  \(\sigma = \mathrm{sgn}(\pi/2 - z)\ln(\sin z)/(6H)\) on one period
  \(0 < z < \pi\); \(\sigma = 0\) is the fold \(z = \pi/2\),
  \(\sigma \to \mp\infty\) the tips \(z \to 0, \pi\); one period of
  \(x_0\) is the whole line \(-\infty < \sigma < \infty\). Unit: length.
  Status: DERIVED.
\item
  \(\mathrm{sgn}\): The sign function: \(+1\) for a positive argument
  and \(-1\) for a negative one. Status: STANDARD.
\item
  \(\pi\): The number pi, the ratio of the circumference of a circle to
  its diameter. Unit: dimensionless. Value: 3.141592654. Status:
  STANDARD.
\item
  \(z\): The angle variable of the hidden direction, \(z = 6Hx_0\); one
  period of \(x_0\) is \(0 < z < \pi\), the fold is \(z = \pi/2\). Unit:
  dimensionless. Status: AUTHOR.
\item
  \(H\): The author's constant (an inverse length) that fixes the period
  \(\pi/(6H)\) of \(x_0\), the warp factor
  \(\mathrm{e}^{-2H\lvert\sigma\rvert}\) of the metric, the cosmological
  constant \(\Lambda = -21H^2\) and the tension
  \(\kappa\lambda_{\mathrm{F}} = 12H\) of the fold; in UM-4 a free
  parameter, bounded from below by the propagation of the gauge fields
  in the bulk. Unit: GeV (1/length). Value: \(H \gtrsim 2.5\) TeV (from
  the invisible width of the Z boson). Status: AUTHOR.
\item
  \(\mathrm{d}\): The differential: \(\mathrm{d}\) in front of a
  quantity is an infinitesimal change of it (as \(\mathrm{d}\sigma\), a
  small step of \(\sigma\)); the quotient of two such changes is a
  derivative; \(\mathrm{d}\) with an exponent, in front of several
  variables, is their volume element, the product of their
  differentials. Status: STANDARD.
\item
  \(x_0\): The hidden space-like coordinate (listed first); the metric
  depends on it through \(z = 6Hx_0\), periodically with the period
  \(\pi/(6H)\). Unit: length. Status: AUTHOR.
\item
  \(\mathrm{e}\): Euler's number, the base of the natural logarithm;
  \(\mathrm{e}^{X}\) is the exponential of \(X\). Unit: dimensionless.
  Value: 2.718281828. Status: STANDARD.
\item
  \(\infty\): Infinity: as a limit of an integral or a sum it means that
  the range has no bound. Status: STANDARD.
\end{itemize}

(the fold is \(\sigma = 0\); the half-period \(0 < z < \pi/2\) is
\(\sigma < 0\), the half-period \(\pi/2 < z < \pi\) is \(\sigma > 0\);
the two ends of the period, the tips \(z \to 0\) and \(z \to \pi\), lie
at \(\sigma \to -\infty\) and \(\sigma \to +\infty\)), the metric of
UM-4 is

\[
g_{MN}\,\mathrm{d}X^M\mathrm{d}X^N = \mathrm{d}\sigma^2 + \mathrm{e}^{-2H\lvert\sigma\rvert}\Big[\mathrm{e}^{2a_4}\big(\mathrm{d}x_1^2
 + \mathrm{d}x_2^2 + \mathrm{d}x_3^2\big) - \mathrm{d}x_4^2 - \mathrm{e}^{-2a_4}b_4^2\big(\mathrm{d}x_5^2 + \mathrm{d}x_6^2 + \mathrm{d}x_7^2\big)\Big] \tag{FD-4}
\] \textbf{Symbols in (FD-4).}

\begin{itemize}
\tightlist
\item
  \(g_{MN}\): The metric tensor of the eight-dimensional spacetime
  (signature \(+,+,+,+,-,-,-,-\) in the order \(x_0, \dots, x_7\)):
  \(g_{MN}\,\mathrm{d}X^M\mathrm{d}X^N\) is the squared interval,
  \(g^{MN}\) its inverse, \(\lvert g\rvert\) the absolute value of its
  determinant (\(\sqrt{\lvert g\rvert}\,\mathrm{d}^8X\) is the proper
  volume element). Unit: dimensionless. Status: AUTHOR.
\item
  \(M\): A coordinate index of eight dimensions (also with a numbered
  copy \(M_1, M_2, \dots\)): it runs over \(x_0, x_1, \dots, x_7\); a
  repeated index is summed. Status: STANDARD.
\item
  \(N\): A coordinate index of eight dimensions (also
  \(N_1, N_2, \dots\)), running over \(x_0, \dots, x_7\); a repeated
  index is summed. Status: STANDARD.
\item
  \(\mathrm{d}\): The differential: \(\mathrm{d}\) in front of a
  quantity is an infinitesimal change of it (as \(\mathrm{d}\sigma\), a
  small step of \(\sigma\)); the quotient of two such changes is a
  derivative; \(\mathrm{d}\) with an exponent, in front of several
  variables, is their volume element, the product of their
  differentials. Status: STANDARD.
\item
  \(X = \{x_0, \dots, x_7\}\): The author's coordinates of the
  eight-dimensional spacetime, \(X = \{x_0, x_1, \dots, x_7\}\) with the
  hidden direction first (Dirac\_Einstein-Lovelock\_4+4\_Nash.nb,
  section DEFCHART); \(X^M\) is the coordinate with the index \(M\) and
  \(\mathrm{d}^8X = \mathrm{d}x_0\,\mathrm{d}x_1\cdots\mathrm{d}x_7\)
  the coordinate volume element. Unit: length (each coordinate). Status:
  AUTHOR.
\item
  \(\sigma\): The proper coordinate of the hidden direction:
  \(\mathrm{d}\sigma = \lvert\cot z\rvert\,\mathrm{d}x_0\),
  \(\sigma = \mathrm{sgn}(\pi/2 - z)\ln(\sin z)/(6H)\) on one period
  \(0 < z < \pi\); \(\sigma = 0\) is the fold \(z = \pi/2\),
  \(\sigma \to \mp\infty\) the tips \(z \to 0, \pi\); one period of
  \(x_0\) is the whole line \(-\infty < \sigma < \infty\). Unit: length.
  Status: DERIVED.
\item
  \(\mathrm{e}\): Euler's number, the base of the natural logarithm;
  \(\mathrm{e}^{X}\) is the exponential of \(X\). Unit: dimensionless.
  Value: 2.718281828. Status: STANDARD.
\item
  \(H\): The author's constant (an inverse length) that fixes the period
  \(\pi/(6H)\) of \(x_0\), the warp factor
  \(\mathrm{e}^{-2H\lvert\sigma\rvert}\) of the metric, the cosmological
  constant \(\Lambda = -21H^2\) and the tension
  \(\kappa\lambda_{\mathrm{F}} = 12H\) of the fold; in UM-4 a free
  parameter, bounded from below by the propagation of the gauge fields
  in the bulk. Unit: GeV (1/length). Value: \(H \gtrsim 2.5\) TeV (from
  the invisible width of the Z boson). Status: AUTHOR.
\item
  \(a_4\): The author's metric function of the time \(x_4\) that scales
  ordinary 3-space: on the fold \(\mathrm{e}^{a_4}\) is the scale factor
  of cosmology and \(a_4'\) the Hubble rate. Unit: dimensionless.
  Status: AUTHOR.
\item
  \(b_4\): The author's metric function of the time \(x_4\) in the
  extra-time components, \(b_4 > 0\), never identically 1 (the author's
  order of 2026-10-24); the scale factor (the factor of the lengths) of
  the extra times is \(\mathrm{e}^{-a_4}b_4\) on the fold and
  \(\mathrm{e}^{-a_4}b_4\,\mathrm{e}^{-H\lvert\sigma\rvert}\) at the
  distance \(\sigma\) from it (the square root of the warp factor times
  \(\mathrm{e}^{-a_4}b_4\)). Unit: dimensionless. Status: AUTHOR.
\item
  \(x_1\): The first coordinate of ordinary 3-space. Unit: length.
  Status: AUTHOR.
\item
  \(x_2\): The second coordinate of ordinary 3-space. Unit: length.
  Status: AUTHOR.
\item
  \(x_3\): The third coordinate of ordinary 3-space. Unit: length.
  Status: AUTHOR.
\item
  \(x_4\): The time; a prime on a function of \(x_4\) alone is
  \(\mathrm{d}/\mathrm{d}x_4\). Unit: length (= time, \(c = 1\)).
  Status: AUTHOR.
\item
  \(x_5\): The first of the three extra time-like coordinates. Unit:
  length. Status: AUTHOR.
\item
  \(x_6\): The second of the three extra time-like coordinates. Unit:
  length. Status: AUTHOR.
\item
  \(x_7\): The third of the three extra time-like coordinates. Unit:
  length. Status: AUTHOR.
\end{itemize}

It depends on \(\sigma\) only through \(\lvert\sigma\rvert\). The
function \(\mathcal{W}(\sigma) = \mathrm{e}^{-2H\lvert\sigma\rvert}\),
which multiplies the squared steps of all seven directions of the fold
(ordinary space, the time and the extra times), is the warp factor of
the metric (its square root \(\mathrm{e}^{-H\lvert\sigma\rvert}\)
multiplies the lengths). The time \(x_4\) is warped like the other
directions (the author's order of 2026-10-25, 13:15 PDT; in UM-3 it was
not), so that every surface \(\sigma = \text{const}\) has the light cone
of the fold. Four properties make the surface \(\sigma = 0\) a fold, and
together they single it out within the period:

\begin{enumerate}
\def\labelenumi{\arabic{enumi}.}
\tightlist
\item
  \emph{It is a mirror.} The reflection \(\sigma \to -\sigma\) (that is
  \(z \to \pi - z\), or \(x_0 \to \pi/(6H) - x_0\)) maps the metric to
  itself, and the fold is the only surface of the period that the
  reflection leaves in place: the two halves of the period are mirror
  images of each other, joined along the fold, as the two halves of a
  folded sheet of paper are joined along the crease.
\item
  \emph{The warp factor is largest on it.} The warp factor
  \((\sin^2z)^{1/6} = \mathrm{e}^{-2H\lvert\sigma\rvert}\) equals 1 on
  the fold and decreases on both sides.
\item
  \emph{The metric has a kink on it.} The metric is continuous across
  the fold, but its first derivative along \(\sigma\) jumps. The
  extrinsic curvature of the surfaces \(\sigma = \text{const}\) changes
  sign across the fold,
\end{enumerate}

\[
K^a{}_b = \frac12\,g^{ac}\,\partial_\sigma g_{cb} = -H\,\mathrm{sgn}(\sigma)\,\delta^a{}_b\quad(\text{all seven directions of the fold}) \tag{FD-5}
\] \textbf{Symbols in (FD-5).}

\begin{itemize}
\tightlist
\item
  \(K^a{}_b\): The extrinsic curvature of a surface of constant
  \(\sigma\) of the hidden direction,
  \(K^a{}_b = \frac12 g^{ac}\partial_\sigma g_{cb}\) for \(a, b\) along
  the surface (FD-5): \(-H\,\mathrm{sgn}(\sigma)\) along all seven
  directions of the fold for the author's form of the metric of UM-4;
  \(K = K^c{}_c\) its trace. Unit: GeV (1/length). Status: DERIVED.
\item
  \(a\): A frame (orthonormal) index running over \(0, 1, \dots, 7\)
  (frame direction \((0)\) along \(\sigma\), \((1), (2), (3)\) along
  ordinary space, \((4)\) along the time, \((5), (6), (7)\) along the
  extra times); raised and lowered with the frame metric
  \(\eta_{ab} = \mathrm{diag}(1, 1, 1, 1, -1, -1, -1, -1)\); on the fold
  (\(a, b \in \{1, \dots, 7\}\)) the frame of the fold; in the tensors
  of the fold written with its coordinates \(x^a\) (its metric
  \(\hat g_{ab}\) and the tensors built from it) a coordinate index of
  the fold, \(a \in \{1, \dots, 7\}\); a repeated index is summed.
  Status: STANDARD.
\item
  \(b\): A frame index (as \(a\)). Status: STANDARD.
\item
  \(c\): A frame index (as \(a\)). Status: STANDARD.
\item
  \(g_{MN}\): The metric tensor of the eight-dimensional spacetime
  (signature \(+,+,+,+,-,-,-,-\) in the order \(x_0, \dots, x_7\)):
  \(g_{MN}\,\mathrm{d}X^M\mathrm{d}X^N\) is the squared interval,
  \(g^{MN}\) its inverse, \(\lvert g\rvert\) the absolute value of its
  determinant (\(\sqrt{\lvert g\rvert}\,\mathrm{d}^8X\) is the proper
  volume element). Unit: dimensionless. Status: AUTHOR.
\item
  \(\partial\): A partial derivative: the derivative with respect to one
  variable at fixed values of the others (\(\partial_\sigma\) with
  respect to \(\sigma\); \(\partial_M\) and \(\partial_\mu\) with
  respect to the coordinate that carries the same index, of the eight
  dimensions or of the fold). Status: STANDARD.
\item
  \(\sigma\): The label \(\sigma\) in \(\partial_\sigma\): the
  derivative with respect to the proper hidden coordinate \(\sigma\).
  Status: STANDARD.
\item
  \(\sigma\): The proper coordinate of the hidden direction:
  \(\mathrm{d}\sigma = \lvert\cot z\rvert\,\mathrm{d}x_0\),
  \(\sigma = \mathrm{sgn}(\pi/2 - z)\ln(\sin z)/(6H)\) on one period
  \(0 < z < \pi\); \(\sigma = 0\) is the fold \(z = \pi/2\),
  \(\sigma \to \mp\infty\) the tips \(z \to 0, \pi\); one period of
  \(x_0\) is the whole line \(-\infty < \sigma < \infty\). Unit: length.
  Status: DERIVED.
\item
  \(H\): The author's constant (an inverse length) that fixes the period
  \(\pi/(6H)\) of \(x_0\), the warp factor
  \(\mathrm{e}^{-2H\lvert\sigma\rvert}\) of the metric, the cosmological
  constant \(\Lambda = -21H^2\) and the tension
  \(\kappa\lambda_{\mathrm{F}} = 12H\) of the fold; in UM-4 a free
  parameter, bounded from below by the propagation of the gauge fields
  in the bulk. Unit: GeV (1/length). Value: \(H \gtrsim 2.5\) TeV (from
  the invisible width of the Z boson). Status: AUTHOR.
\item
  \(\mathrm{sgn}\): The sign function: \(+1\) for a positive argument
  and \(-1\) for a negative one. Status: STANDARD.
\item
  \(\delta\): The Kronecker delta: with two indices \(\delta^M_N = 1\)
  if \(M = N\) and 0 otherwise (likewise \(\delta_{X\mathrm{L}}\) is 1
  if the label \(X\) of a half of dirac16complex is L and 0 otherwise);
  with two index strings of equal length the generalised Kronecker delta
  \(\delta^{M_1M_2\cdots}_{N_1N_2\cdots}\), the determinant of the
  square matrix of the Kronecker deltas \(\delta^{M_a}_{N_b}\) (\(+1\)
  or \(-1\) if the upper indices are an even or odd permutation of the
  lower ones, 0 otherwise). Status: STANDARD.
\item
  \(\eta_{ab}\): The flat metric of the frame,
  \(\eta = \mathrm{diag}(1, 1, 1, 1, -1, -1, -1, -1)\) (eta4488 of
  4+4.nb); \(\eta^{ab}\) is its inverse (the same matrix). Unit:
  dimensionless. Status: AUTHOR.
\item
  \(\hat g_{ab}\): The metric of the seven-dimensional fold (FD-2) in
  its coordinates \(x_1, \dots, x_7\) (or in its frame), \(\hat g^{ab}\)
  its inverse; \(\hat h_{\mu\nu}\) is its part along the four observable
  directions. Unit: dimensionless. Status: AUTHOR.
\item
  \(\mathrm{diag}\): A diagonal matrix with the listed entries on its
  diagonal (in the order of the coordinates \(x_0, x_1, \dots, x_7\) for
  a tensor of eight dimensions, \(x_1, \dots, x_7\) for a tensor of the
  fold, \(x_1, \dots, x_4\) for a four-dimensional tensor). Status:
  STANDARD.
\item
  \(x\): The coordinates of the fold: \(x_1, x_2, x_3\) of ordinary
  space, \(x_4\) the time, \(x_5, x_6, x_7\) the extra times; a field of
  the fold is a function of the four observable coordinates
  \(x = (x_1, x_2, x_3, x_4)\), and
  \(\mathrm{d}^4x = \mathrm{d}x_1\,\mathrm{d}x_2\,\mathrm{d}x_3\,\mathrm{d}x_4\).
  Unit: length (each coordinate). Status: AUTHOR.
\end{itemize}

(\(a, b\) frame indices along the fold); it jumps by \(-2H\) in every
direction of the fold. The field equations of gravity of UM-4,
\(\mathcal{E}_{(1)}{}^M{}_N + \Lambda\delta^M_N = \kappa\mathcal{T}^M{}_N\)
(Einstein gravity: the author's Einstein-Lovelock equations with the
weights of the Gauss-Bonnet and cubic terms set to zero;
\(\mathcal{E}_{(1)}\) the Einstein tensor, \(\Lambda = -21H^2\) the
cosmological constant, \(\kappa = 8\pi G_8\) the gravitational coupling
of eight dimensions, \(\mathcal{T}^M{}_N\) the energy-momentum tensor,
\(\delta^M_N\) the Kronecker delta), are linear in the second
\(\sigma\)-derivatives of the metric, so the jump puts a Dirac delta
function \(\delta_{\mathrm{D}}(\sigma)\) into them: the metric solves
them only if the fold carries a thin source \(\mathcal{S}^a{}_b\), an
energy-momentum tensor concentrated on \(\sigma = 0\), given by the
integral of the equations across the fold (the junction condition of
Israel),

\[
\kappa\,\mathcal{S}^a{}_b = \int_{0^-}^{0^+}\mathrm{d}\sigma\,\mathcal{E}_{(1)}{}^a{}_b
 = -\Big(\big(K^a{}_b\big\vert_{0^+} - K^a{}_b\big\vert_{0^-}\big) - \delta^a{}_b\big(K\big\vert_{0^+} - K\big\vert_{0^-}\big)\Big) \tag{FD-11}
\] \textbf{Symbols in (FD-11).}

\begin{itemize}
\tightlist
\item
  \(\kappa\): The eight-dimensional gravitational coupling,
  \(\kappa = 8\pi G_8\). Unit: GeV\(^{-6}\). Status: AUTHOR.
\item
  \(\mathcal{S}^a{}_b\): The thin source of the fold: the
  energy-momentum tensor concentrated on the fold, integrated across it,
  \(\kappa\mathcal{S}^a{}_b = \int_{0^-}^{0^+}\mathrm{d}\sigma\,\mathcal{E}_{(1)}{}^a{}_b\)
  (the junction condition of Israel); \(\mathcal{S}^4{}_4\) is minus its
  energy density, \(\mathcal{S}^1{}_1\) and \(\mathcal{S}^5{}_5\) its
  pressures along ordinary space and along the extra times (frame
  components along the fold, eight-dimensional normalisation: an energy
  per unit volume of the six other directions of the fold). Unit:
  GeV\(^7\). Value: \(\kappa\mathcal{S}^a{}_b = -12H\delta^a{}_b\) for
  every \(a_4\), \(b_4\) of the author's form. Status: DERIVED.
\item
  \(a\): A frame (orthonormal) index running over \(0, 1, \dots, 7\)
  (frame direction \((0)\) along \(\sigma\), \((1), (2), (3)\) along
  ordinary space, \((4)\) along the time, \((5), (6), (7)\) along the
  extra times); raised and lowered with the frame metric
  \(\eta_{ab} = \mathrm{diag}(1, 1, 1, 1, -1, -1, -1, -1)\); on the fold
  (\(a, b \in \{1, \dots, 7\}\)) the frame of the fold; in the tensors
  of the fold written with its coordinates \(x^a\) (its metric
  \(\hat g_{ab}\) and the tensors built from it) a coordinate index of
  the fold, \(a \in \{1, \dots, 7\}\); a repeated index is summed.
  Status: STANDARD.
\item
  \(b\): A frame index (as \(a\)). Status: STANDARD.
\item
  \(\mathrm{d}\): The differential: \(\mathrm{d}\) in front of a
  quantity is an infinitesimal change of it (as \(\mathrm{d}\sigma\), a
  small step of \(\sigma\)); the quotient of two such changes is a
  derivative; \(\mathrm{d}\) with an exponent, in front of several
  variables, is their volume element, the product of their
  differentials. Status: STANDARD.
\item
  \(\sigma\): The proper coordinate of the hidden direction:
  \(\mathrm{d}\sigma = \lvert\cot z\rvert\,\mathrm{d}x_0\),
  \(\sigma = \mathrm{sgn}(\pi/2 - z)\ln(\sin z)/(6H)\) on one period
  \(0 < z < \pi\); \(\sigma = 0\) is the fold \(z = \pi/2\),
  \(\sigma \to \mp\infty\) the tips \(z \to 0, \pi\); one period of
  \(x_0\) is the whole line \(-\infty < \sigma < \infty\). Unit: length.
  Status: DERIVED.
\item
  \(\mathcal{E}_{(p)}{}^M{}_N\): The Lovelock tensor of order \(p\),
  \(\mathcal{E}_{(p)}{}^M{}_N = -2^{-(p+1)}\,\delta^{MP_1\cdots P_{2p}}_{NQ_1\cdots Q_{2p}}\,R^{Q_1Q_2}{}_{P_1P_2}\cdots R^{Q_{2p-1}Q_{2p}}{}_{P_{2p-1}P_{2p}}\)
  (with both indices down it is \(1/\sqrt{\lvert g\rvert}\) times the
  variation of \(\sqrt{\lvert g\rvert}\,\mathcal{L}_{(p)}\) with respect
  to \(g^{MN}\); \(\mathcal{E}_{(1)}\) the Einstein tensor,
  \(\mathcal{E}_{(2)}\) the Gauss-Bonnet tensor, \(\mathcal{E}_{(3)}\)
  the cubic Lovelock tensor); with frame indices its frame components.
  Unit: 1/length\(^{2p}\). Status: STANDARD.
\item
  \(p\): The Lovelock order \(p = 1, 2, 3\) (Einstein, Gauss-Bonnet,
  cubic), as an index and as a number (UM-4 keeps only \(p = 1\)); in
  the Kohn-Sham equations also a second label of a Kohn-Sham orbital,
  with the Kronecker delta of the two labels. Status: STANDARD.
\item
  \(K^a{}_b\): The extrinsic curvature of a surface of constant
  \(\sigma\) of the hidden direction,
  \(K^a{}_b = \frac12 g^{ac}\partial_\sigma g_{cb}\) for \(a, b\) along
  the surface (FD-5): \(-H\,\mathrm{sgn}(\sigma)\) along all seven
  directions of the fold for the author's form of the metric of UM-4;
  \(K = K^c{}_c\) its trace. Unit: GeV (1/length). Status: DERIVED.
\item
  \(\delta\): The Kronecker delta: with two indices \(\delta^M_N = 1\)
  if \(M = N\) and 0 otherwise (likewise \(\delta_{X\mathrm{L}}\) is 1
  if the label \(X\) of a half of dirac16complex is L and 0 otherwise);
  with two index strings of equal length the generalised Kronecker delta
  \(\delta^{M_1M_2\cdots}_{N_1N_2\cdots}\), the determinant of the
  square matrix of the Kronecker deltas \(\delta^{M_a}_{N_b}\) (\(+1\)
  or \(-1\) if the upper indices are an even or odd permutation of the
  lower ones, 0 otherwise). Status: STANDARD.
\item
  \(G_8\): The eight-dimensional Newton constant
  (\(\kappa = 8\pi G_8\)). Unit: GeV\(^{-6}\). Status: AUTHOR.
\item
  \(R^{MN}{}_{PQ},\ R_{MN},\ R\): The curvature of the metric
  \(g_{MN}\): with four indices the Riemann tensor \(R^M{}_{NPQ}\) (in
  the convention in which the commutator of two covariant derivatives,
  acting on a vector, is the Riemann tensor, with its last two indices
  those of the two derivatives, contracted with the vector), with two
  indices the Ricci tensor \(R_{MN} = R^P{}_{MPN}\), without indices the
  scalar curvature \(R = g^{MN}R_{MN}\); with frame indices
  \(R^{(a)(b)}{}_{(c)(d)}\) its frame components,
  \(R^{(a)(b)}{}_{(a)(b)}\) (no sum) the sectional curvature of the
  frame plane \((a, b)\). Unit: 1/length\(^2\). Status: STANDARD.
\item
  \(\mathcal{L}_{(p)}\): The Lovelock density of order \(p\),
  \(\mathcal{L}_{(p)} = 2^{-p}\,\delta^{P_1\cdots P_{2p}}_{Q_1\cdots Q_{2p}}\,R^{Q_1Q_2}{}_{P_1P_2}\cdots R^{Q_{2p-1}Q_{2p}}{}_{P_{2p-1}P_{2p}}\),
  a scalar built from \(p\) curvature tensors: \(\mathcal{L}_{(1)} = R\)
  (Einstein-Hilbert), \(\mathcal{L}_{(2)}\) the Gauss-Bonnet density,
  \(\mathcal{L}_{(3)}\) the cubic Lovelock density. Unit:
  1/length\(^{2p}\). Status: STANDARD.
\end{itemize}

(\(K = K^c{}_c\) the trace of the extrinsic curvature), with
\(\mathcal{S}^4{}_4 = -(\text{its energy density})\),
\(\mathcal{S}^1{}_1 = \mathcal{S}^2{}_2 = \mathcal{S}^3{}_3\) its
pressure along ordinary space and
\(\mathcal{S}^5{}_5 = \mathcal{S}^6{}_6 = \mathcal{S}^7{}_7\) its
pressure along the extra times; \(\mathcal{S}\) is \(\mathcal{T}\)
integrated across the fold, an energy per unit volume of the six other
directions of the fold (unit GeV\(^7\)). For the static block (\(a_4\)
and \(b_4\) constant) and, to the first order in the curvature of the
fold divided by \(H^2\), for every fold

\[
\begin{aligned}
&\text{static block:}\quad \kappa\mathcal{S}^a{}_b = -12H\,\delta^a{}_b\quad\text{(a pure tension)},\qquad \kappa\lambda_{\mathrm{F}} = 12H,\\
&\text{any fold:}\quad \kappa\mathcal{S}^a{}_b = -12H\,\delta^a{}_b + \frac{2}{5H}\,\hat G^a{}_b + \dots
\end{aligned} \tag{FD-12}
\] \textbf{Symbols in (FD-12).}

\begin{itemize}
\tightlist
\item
  \(\kappa\): The eight-dimensional gravitational coupling,
  \(\kappa = 8\pi G_8\). Unit: GeV\(^{-6}\). Status: AUTHOR.
\item
  \(\mathcal{S}^a{}_b\): The thin source of the fold: the
  energy-momentum tensor concentrated on the fold, integrated across it,
  \(\kappa\mathcal{S}^a{}_b = \int_{0^-}^{0^+}\mathrm{d}\sigma\,\mathcal{E}_{(1)}{}^a{}_b\)
  (the junction condition of Israel); \(\mathcal{S}^4{}_4\) is minus its
  energy density, \(\mathcal{S}^1{}_1\) and \(\mathcal{S}^5{}_5\) its
  pressures along ordinary space and along the extra times (frame
  components along the fold, eight-dimensional normalisation: an energy
  per unit volume of the six other directions of the fold). Unit:
  GeV\(^7\). Value: \(\kappa\mathcal{S}^a{}_b = -12H\delta^a{}_b\) for
  every \(a_4\), \(b_4\) of the author's form. Status: DERIVED.
\item
  \(a\): A frame (orthonormal) index running over \(0, 1, \dots, 7\)
  (frame direction \((0)\) along \(\sigma\), \((1), (2), (3)\) along
  ordinary space, \((4)\) along the time, \((5), (6), (7)\) along the
  extra times); raised and lowered with the frame metric
  \(\eta_{ab} = \mathrm{diag}(1, 1, 1, 1, -1, -1, -1, -1)\); on the fold
  (\(a, b \in \{1, \dots, 7\}\)) the frame of the fold; in the tensors
  of the fold written with its coordinates \(x^a\) (its metric
  \(\hat g_{ab}\) and the tensors built from it) a coordinate index of
  the fold, \(a \in \{1, \dots, 7\}\); a repeated index is summed.
  Status: STANDARD.
\item
  \(b\): A frame index (as \(a\)). Status: STANDARD.
\item
  \(H\): The author's constant (an inverse length) that fixes the period
  \(\pi/(6H)\) of \(x_0\), the warp factor
  \(\mathrm{e}^{-2H\lvert\sigma\rvert}\) of the metric, the cosmological
  constant \(\Lambda = -21H^2\) and the tension
  \(\kappa\lambda_{\mathrm{F}} = 12H\) of the fold; in UM-4 a free
  parameter, bounded from below by the propagation of the gauge fields
  in the bulk. Unit: GeV (1/length). Value: \(H \gtrsim 2.5\) TeV (from
  the invisible width of the Z boson). Status: AUTHOR.
\item
  \(\delta\): The Kronecker delta: with two indices \(\delta^M_N = 1\)
  if \(M = N\) and 0 otherwise (likewise \(\delta_{X\mathrm{L}}\) is 1
  if the label \(X\) of a half of dirac16complex is L and 0 otherwise);
  with two index strings of equal length the generalised Kronecker delta
  \(\delta^{M_1M_2\cdots}_{N_1N_2\cdots}\), the determinant of the
  square matrix of the Kronecker deltas \(\delta^{M_a}_{N_b}\) (\(+1\)
  or \(-1\) if the upper indices are an even or odd permutation of the
  lower ones, 0 otherwise). Status: STANDARD.
\item
  \(\lambda_{\mathrm{F}}\): The tension that the static block requires
  of the fold, \(\kappa\lambda_{\mathrm{F}} = 12H\): the total vacuum
  energy of the seven-dimensional fold, made of the tension of the fold
  fluid and the vacuum energy of the fields of the fold (both have the
  pressure minus their energy density along ordinary space and along the
  extra times), per unit volume of the six other directions of the fold
  (eight-dimensional normalisation); in the four-dimensional
  normalisation it is \(\lambda_{\mathrm{F}}V_*\), and the excess of the
  total over \(\lambda_{\mathrm{F}}V_*\) is the dark energy
  \(\rho_\Lambda\). Unit: GeV\(^7\). Value: \(12H/\kappa\);
  \(\lambda_{\mathrm{F}}V_* = 30H^2M_{\mathrm{red}}^2 = 1.1\times10^{45}\)
  GeV\(^4\) for \(H = 2.5\) TeV. Status: DERIVED.
\item
  \(\hat G^\mu{}_\nu\): The Einstein tensor of the metric
  \(\hat g_{ab}\) of the fold (FD-2), with frame indices along its seven
  directions; with four-dimensional indices, \(\hat G^\mu{}_\nu\), its
  components along the four observable directions, which for frozen
  extra times are those of the Einstein tensor of the four-dimensional
  metric \(\hat h_{\mu\nu}\). Unit: 1/length\(^2\). Status: STANDARD.
\item
  \(M_{\mathrm{red}}\): The reduced Planck mass,
  \(M_{\mathrm{red}} = (8\pi G_{\mathrm{N}})^{-1/2}\). Unit: GeV. Value:
  \(2.435\times10^{18}\) GeV. Status: STANDARD.
\item
  \(V_*\): The physical volume of the extra times \(x_5, x_6, x_7\) on
  the fold in the late universe,
  \(V_* = V_{\mathrm{t}}(\mathrm{e}^{-a_4}b_4)^3\) (with
  \(V_{\mathrm{t}} = \int\mathrm{d}x_5\,\mathrm{d}x_6\,\mathrm{d}x_7\)
  the coordinate volume of the extra times) at a time when
  \(\mathrm{e}^{-a_4}b_4\) is constant; it converts the
  eight-dimensional normalisation into the four-dimensional one. Unit:
  length\(^3\) (GeV\(^{-3}\)). Status: ASSUMED.
\item
  \(\hat h_{\mu\nu}\): The four-dimensional metric of the fold,
  \(\hat h_{\mu\nu} = \mathrm{diag}(\mathrm{e}^{2a_4}, \mathrm{e}^{2a_4}, \mathrm{e}^{2a_4}, -1)\)
  in the coordinates \(x_1, \dots, x_4\) (the induced metric at
  \(\sigma = 0\) without the extra times); \(\hat h^{\mu\nu}\) its
  inverse, \(\hat h\) its determinant,
  \(\sqrt{-\hat h} = \mathrm{e}^{3a_4}\); \(\hat h^{(1)}_{\mu\nu}\) its
  first-order cosmological perturbation. Unit: dimensionless. Status:
  DERIVED.
\item
  \(\rho_\Lambda\): The vacuum energy density of the fold above its
  static tension, the dark energy of UM-4:
  \(\rho_\Lambda = \rho_{\mathrm{ff}} + \rho_{\mathrm{vac}} - \lambda_{\mathrm{F}}V_*\)
  (four-dimensional normalisation). Unit: GeV\(^4\). Value:
  \(0.685\,\rho_{\mathrm{c}}^{(0)} = 2.5\times10^{-47}\) GeV\(^4\)
  (Planck 2018). Status: DERIVED.
\item
  \(G_{\mathrm{N}}\): Newton's constant of four-dimensional gravity (the
  gravity of our world, on the fold),
  \(G_{\mathrm{N}} = 1/M_{\mathrm{Pl}}^2\) with the Planck mass
  \(M_{\mathrm{Pl}} = 1.221\times10^{19}\) GeV. Unit: GeV\(^{-2}\).
  Value: \(6.708\times10^{-39}\) GeV\(^{-2}\). Status: STANDARD.
\item
  \(V_{\mathrm{t}}\): The coordinate volume of the (compact) extra
  times,
  \(V_{\mathrm{t}} = \int\mathrm{d}x_5\,\mathrm{d}x_6\,\mathrm{d}x_7\).
  Unit: length\(^3\). Status: ASSUMED.
\item
  \(\rho_{\mathrm{c}}\): The critical density, the energy density of a
  spatially flat fold at the expansion rate \(H_4\);
  \(\rho_{\mathrm{c}}^{(0)}\) its value today. Unit: GeV\(^4\). Value:
  \(\rho_{\mathrm{c}}^{(0)} = 1.0537\times10^{-5}\,h_{100}^2\)
  GeV/cm\(^3\). Status: STANDARD.
\item
  \(\rho_{\mathrm{ff}}\): The energy density of the fold fluid. Unit:
  GeV\(^4\). Value:
  \(\lambda_{\mathrm{F}}V_* + \rho_\Lambda - \rho_{\mathrm{vac}}\):
  together with the vacuum energy \(\rho_{\mathrm{vac}}\) of the fields
  of the fold it makes up the tension \(\lambda_{\mathrm{F}}V_*\) that
  the static block requires plus the dark energy \(\rho_\Lambda\).
  Status: ASSUMED (the fluid), DERIVED (the value).
\item
  \(\rho_{\mathrm{vac}}\): The vacuum energy of the fields of the fold:
  the minimum of their potential (for the Standard Model the minimum of
  the Higgs potential with its loop corrections), a tension of the
  seven-dimensional fold. Unit: GeV\(^4\). Status: STANDARD.
\item
  \(H_4\): The expansion rate (Hubble rate) of the fold, \(H_4 = a_4'\),
  as a function of the time or of the redshift; \(H_4^{(0)}\) its value
  today, the Hubble constant. Unit: GeV (1/time). Value:
  \(H_4^{(0)} = 100\,h_{100}\) km/s/Mpc with \(h_{100} = 0.674\) (Planck
  2018): 67.4 km/s/Mpc \(= 1.44\times10^{-42}\) GeV. Status: STANDARD.
\item
  \(M_{\mathrm{Pl}}\): The Planck mass,
  \(M_{\mathrm{Pl}} = G_{\mathrm{N}}^{-1/2}\). Unit: GeV. Value:
  \(1.221\times10^{19}\) GeV. Status: STANDARD.
\item
  \(h_{100}\): The present expansion rate of the fold in units of 100
  km/s/Mpc. Unit: dimensionless. Value: 0.674 (Planck 2018). Status:
  STANDARD.
\end{itemize}

(\(\hat G^a{}_b\) the Einstein tensor of the metric FD-2 of the fold):
the tension plus matter that obeys the Einstein equations of seven
dimensions, \(\hat G^a{}_b = \kappa_7\mathsf{t}^a{}_b\) with
\(\kappa_7 = 5\kappa H/2\) (\(\mathsf{t}^a{}_b\) the matter of the
fold), which for frozen extra times are the Friedmann equations of
general relativity. The fold is therefore not only a place but a
physical object: a thin layer that carries energy and momentum (a
brane). 4. \emph{The author's coordinate \(x_0\) degenerates on it.} The
metric component \(g_{x_0x_0} = \cot^2(6Hx_0)\) vanishes on the fold and
nowhere else in the period; near the fold the proper distance grows with
the square of the coordinate distance,

\[
\lvert\sigma\rvert = 3H\Big(x_0 - \frac{\pi}{12H}\Big)^2\Big[1 + 6H^2\Big(x_0 - \frac{\pi}{12H}\Big)^2 + \dots\Big] \tag{FD-6}
\] \textbf{Symbols in (FD-6).}

\begin{itemize}
\tightlist
\item
  \(\sigma\): The proper coordinate of the hidden direction:
  \(\mathrm{d}\sigma = \lvert\cot z\rvert\,\mathrm{d}x_0\),
  \(\sigma = \mathrm{sgn}(\pi/2 - z)\ln(\sin z)/(6H)\) on one period
  \(0 < z < \pi\); \(\sigma = 0\) is the fold \(z = \pi/2\),
  \(\sigma \to \mp\infty\) the tips \(z \to 0, \pi\); one period of
  \(x_0\) is the whole line \(-\infty < \sigma < \infty\). Unit: length.
  Status: DERIVED.
\item
  \(H\): The author's constant (an inverse length) that fixes the period
  \(\pi/(6H)\) of \(x_0\), the warp factor
  \(\mathrm{e}^{-2H\lvert\sigma\rvert}\) of the metric, the cosmological
  constant \(\Lambda = -21H^2\) and the tension
  \(\kappa\lambda_{\mathrm{F}} = 12H\) of the fold; in UM-4 a free
  parameter, bounded from below by the propagation of the gauge fields
  in the bulk. Unit: GeV (1/length). Value: \(H \gtrsim 2.5\) TeV (from
  the invisible width of the Z boson). Status: AUTHOR.
\item
  \(x_0\): The hidden space-like coordinate (listed first); the metric
  depends on it through \(z = 6Hx_0\), periodically with the period
  \(\pi/(6H)\). Unit: length. Status: AUTHOR.
\item
  \(\pi\): The number pi, the ratio of the circumference of a circle to
  its diameter. Unit: dimensionless. Value: 3.141592654. Status:
  STANDARD.
\end{itemize}

so the two coordinate values on either side of \(x_0 = \pi/(12H)\) at
the same coordinate distance lie at the same proper distance from the
fold: the coordinate line of \(x_0\) is folded onto the line of proper
distances (a fold map, example 5).

\textbf{What lives on the fold.} In UM-4 the Higgs doublet (with the
author's ad hoc vacuum, a time-dependent choice of its vacuum value),
three generations of quarks and leptons, three right-handed neutrinos
and the fold fluid (its tension and the pressure along the extra times
that keeps their volume fixed) live on the fold: they are fields at
\(\sigma = 0\), functions of \(x_1, \dots, x_4\) (no field depends on
the extra times). The gauge fields of
\(SU(3)_C \times SU(2)_L \times U(1)_Y\) live in the bulk; their
massless modes are bound to the fold, so that the photon and the gluons
of the four-dimensional world are waves of the bulk that the fold holds,
and photons that leave the fold travel through the bulk at the speed of
light. dirac16complex (the author's spinor field of eight dimensions)
and gravity live in the bulk \(\sigma \ne 0\) as well as on the fold;
the states \(\chi\) of dirac16complex that the fold holds (in UM-4
long-lived resonances) are the part of it near the fold, and the portals
(couplings between dirac16complex and the fields of the fold) act at
\(\sigma = 0\). Integrated over the extra times, the fold is the
four-dimensional universe that we observe, with the metric
\(\mathrm{diag}(\mathrm{e}^{2a_4}, \mathrm{e}^{2a_4}, \mathrm{e}^{2a_4}, -1)\)
in \(x_1, \dots, x_4\). In UM-4 files ``the fold'' means this surface;
the examples below use the word for other surfaces with the same
properties. The word is not used in its other senses: a set of equal
energy levels is called degenerate, and a growth by the factor
\(\mathrm{e}\) is called an e-folding.

\textbf{Examples of folds that are meaningful for UM-4.}

\begin{enumerate}
\def\labelenumi{\arabic{enumi}.}
\tightlist
\item
  \emph{Our fold.} The fold \(z = \pi/2\) of the period that contains
  our world (FD-1): the seven-dimensional world-volume of the quarks,
  leptons and the Higgs field, and after the integration over the extra
  times the four-dimensional universe that we observe.
\item
  \emph{The copies of the fold.} Because the metric repeats with the
  period \(\pi/(6H)\), every surface
  \(x_0 = \pi/(12H) + (\text{an integer multiple of } \pi/(6H))\) is a
  copy of the fold, one per period. Neighbouring copies are separated by
  a tip (\(z\) a multiple of \(\pi\)), where the warp factor
  \((\sin^2z)^{1/6}\) vanishes and \(g_{x_0x_0} = \cot^2z\) diverges,
  and which lies at an infinite proper distance from both: in proper
  distance every copy has its own complete line
  \(-\infty < \sigma < \infty\). A compact direction of length \(L\)
  would give every field a tower of particles with the masses
  \(2\pi n/L\), \(n = 1, 2, \dots\) (a Kaluza-Klein tower); because the
  copies are infinitely far apart, the period \(\pi/(6H)\) of \(x_0\)
  produces no such tower with the spacing \(12H\): every field of the
  bulk has instead a continuum of modes that starts at zero mass, and
  the fold holds its massless mode.
\item
  \emph{A folded sheet of paper.} The crease of a folded sheet is a fold
  in this sense: the two halves are mirror images (property 1) joined
  along the crease, the sheet is continuous there and its direction
  changes abruptly (property 3). In UM-4 the direction across the crease
  is the hidden direction \(\sigma\), the seven directions along the
  crease are \(x_1, \dots, x_7\), and the crease is the fold.
\item
  \emph{A folded function.} The warp factor
  \(\mathrm{e}^{-2H\lvert\sigma\rvert}\) is the simplest folded
  function: continuous, mirror symmetric, largest at \(\sigma = 0\),
  with a jump of its derivative,
\end{enumerate}

\[
\frac{\mathrm{d}}{\mathrm{d}\sigma}\mathrm{e}^{-2H\lvert\sigma\rvert} = -2H\,\mathrm{sgn}(\sigma)\,\mathrm{e}^{-2H\lvert\sigma\rvert},\qquad
\frac{\mathrm{d}^2}{\mathrm{d}\sigma^2}\mathrm{e}^{-2H\lvert\sigma\rvert} = 4H^2\,\mathrm{e}^{-2H\lvert\sigma\rvert} - 4H\,\delta_{\mathrm{D}}(\sigma) \tag{FD-7}
\] \textbf{Symbols in (FD-7).}

\begin{itemize}
\tightlist
\item
  \(\mathrm{d}\): The differential: \(\mathrm{d}\) in front of a
  quantity is an infinitesimal change of it (as \(\mathrm{d}\sigma\), a
  small step of \(\sigma\)); the quotient of two such changes is a
  derivative; \(\mathrm{d}\) with an exponent, in front of several
  variables, is their volume element, the product of their
  differentials. Status: STANDARD.
\item
  \(\sigma\): The proper coordinate of the hidden direction:
  \(\mathrm{d}\sigma = \lvert\cot z\rvert\,\mathrm{d}x_0\),
  \(\sigma = \mathrm{sgn}(\pi/2 - z)\ln(\sin z)/(6H)\) on one period
  \(0 < z < \pi\); \(\sigma = 0\) is the fold \(z = \pi/2\),
  \(\sigma \to \mp\infty\) the tips \(z \to 0, \pi\); one period of
  \(x_0\) is the whole line \(-\infty < \sigma < \infty\). Unit: length.
  Status: DERIVED.
\item
  \(\mathrm{e}\): Euler's number, the base of the natural logarithm;
  \(\mathrm{e}^{X}\) is the exponential of \(X\). Unit: dimensionless.
  Value: 2.718281828. Status: STANDARD.
\item
  \(H\): The author's constant (an inverse length) that fixes the period
  \(\pi/(6H)\) of \(x_0\), the warp factor
  \(\mathrm{e}^{-2H\lvert\sigma\rvert}\) of the metric, the cosmological
  constant \(\Lambda = -21H^2\) and the tension
  \(\kappa\lambda_{\mathrm{F}} = 12H\) of the fold; in UM-4 a free
  parameter, bounded from below by the propagation of the gauge fields
  in the bulk. Unit: GeV (1/length). Value: \(H \gtrsim 2.5\) TeV (from
  the invisible width of the Z boson). Status: AUTHOR.
\item
  \(\mathrm{sgn}\): The sign function: \(+1\) for a positive argument
  and \(-1\) for a negative one. Status: STANDARD.
\item
  \(\delta_{\mathrm{D}}\): The Dirac delta function of \(\sigma\): zero
  for \(\sigma \ne 0\), with the integral 1 over any range that contains
  \(\sigma = 0\); the integral of \(\delta_{\mathrm{D}}(\sigma)\) times
  a smooth function of \(\sigma\) is the value of that function at
  \(\sigma = 0\). Unit: 1/length (for \(\delta_{\mathrm{D}}(\sigma)\)).
  Status: STANDARD.
\end{itemize}

whose delta function is the origin of the thin source of property 3.
Every metric function of FD-4 depends on \(\sigma\) through
\(\lvert\sigma\rvert\) and is folded at \(\sigma = 0\) in this sense. 5.
\emph{A fold map.} In the mathematics of smooth maps the fold (Whitney's
fold singularity) is the map \(u \mapsto u^2\) of a line onto a
half-line: it folds the line at \(u = 0\), two points \(\pm u\) land on
the same point, and the derivative of the map vanishes on the fold. Near
the fold the map from the author's coordinate \(x_0\) to the distance
\(\lvert\sigma\rvert\) from the fold is of this kind (FD-6), and
\(g_{x_0x_0} = \cot^2z\), the square of the derivative
\(\mathrm{d}\sigma/\mathrm{d}x_0\), vanishes there. 6. \emph{The brane
of Randall and Sundrum.} The five-dimensional metric

\[
\mathrm{d}y_{\mathrm{RS}}^2 + \mathrm{e}^{-2k_{\mathrm{RS}}\lvert y_{\mathrm{RS}}\rvert}\big(\mathrm{d}x_1^2 + \mathrm{d}x_2^2 + \mathrm{d}x_3^2
 - \mathrm{d}x_4^2\big),\qquad \lambda_{\mathrm{RS}} = \frac{6k_{\mathrm{RS}}}{\kappa_5},\qquad \Lambda_{\mathrm{RS}} = -\frac{6k_{\mathrm{RS}}^2}{\kappa_5} \tag{FD-8}
\] \textbf{Symbols in (FD-8).}

\begin{itemize}
\tightlist
\item
  \(\mathrm{d}\): The differential: \(\mathrm{d}\) in front of a
  quantity is an infinitesimal change of it (as \(\mathrm{d}\sigma\), a
  small step of \(\sigma\)); the quotient of two such changes is a
  derivative; \(\mathrm{d}\) with an exponent, in front of several
  variables, is their volume element, the product of their
  differentials. Status: STANDARD.
\item
  \(y_{\mathrm{RS}}\): The proper coordinate across the brane of the
  five-dimensional model of Randall and Sundrum (FD-8, an example, not
  UM-4). Unit: length. Status: STANDARD.
\item
  \(\mathrm{e}\): Euler's number, the base of the natural logarithm;
  \(\mathrm{e}^{X}\) is the exponential of \(X\). Unit: dimensionless.
  Value: 2.718281828. Status: STANDARD.
\item
  \(k_{\mathrm{RS}}\): The inverse curvature length of the bulk of the
  model of Randall and Sundrum: the warp factor of the brane is
  \(\mathrm{e}^{-2k_{\mathrm{RS}}\lvert y_{\mathrm{RS}}\rvert}\) (FD-8,
  an example, not UM-4). Unit: GeV (1/length). Status: STANDARD.
\item
  \(x_1\): The first coordinate of ordinary 3-space. Unit: length.
  Status: AUTHOR.
\item
  \(x_2\): The second coordinate of ordinary 3-space. Unit: length.
  Status: AUTHOR.
\item
  \(x_3\): The third coordinate of ordinary 3-space. Unit: length.
  Status: AUTHOR.
\item
  \(x_4\): The time; a prime on a function of \(x_4\) alone is
  \(\mathrm{d}/\mathrm{d}x_4\). Unit: length (= time, \(c = 1\)).
  Status: AUTHOR.
\item
  \(\lambda_{\mathrm{RS}}\): The tension (energy per unit volume of
  ordinary space) of the brane of the model of Randall and Sundrum,
  \(\lambda_{\mathrm{RS}} = 6k_{\mathrm{RS}}/\kappa_5\) (FD-8, an
  example, not UM-4). Unit: GeV\(^4\). Status: STANDARD.
\item
  \(\kappa_5\): The gravitational coupling of five dimensions in the
  model of Randall and Sundrum, \(8\pi\) times the Newton constant of
  five dimensions (FD-8, an example, not UM-4). Unit: GeV\(^{-3}\).
  Status: STANDARD.
\item
  \(\Lambda_{\mathrm{RS}}\): The vacuum energy density of the bulk of
  the model of Randall and Sundrum,
  \(\Lambda_{\mathrm{RS}} = -6k_{\mathrm{RS}}^2/\kappa_5\) (FD-8, an
  example, not UM-4). Unit: GeV\(^5\). Status: STANDARD.
\end{itemize}

has a fold at \(y_{\mathrm{RS}} = 0\) with the properties 1 to 3 (its
coordinate \(y_{\mathrm{RS}}\) is already the proper distance, so
property 4 does not arise); Einstein's equations require there a brane
of the tension \(\lambda_{\mathrm{RS}}\) (Israel's junction condition)
and in the bulk the vacuum energy \(\Lambda_{\mathrm{RS}}\), and the
bulk is otherwise empty. The fold of UM-4 is of the same kind, with
every direction of the fold warped by the same factor as in the model of
Randall and Sundrum: its static block is a space of constant curvature
\(-H^2\) in all eight directions, a vacuum of Einstein gravity with
\(\Lambda = -21H^2\), with a fold of pure tension (FD-12), and the
graviton is bound to the fold as there. It differs in that the fold has
seven dimensions, three of them extra times, in that the space is not
cut at the brane (Randall and Sundrum identify the two sides of their
brane, one physical side; the fold of UM-4 has two distinct sides,
example 7), and in that the gauge fields live in the bulk as well. 7.
\emph{A particle held by a thin layer.} A particle held by an attractive
thin layer has a wave function that falls off exponentially with the
distance on both sides, and the layer makes its derivative jump (for the
first-order Dirac equation, the wave function itself). In UM-4 the
states \(\chi\) of dirac16complex near the fold are of this kind only
near the fold: at rest and without their self-interaction (the field is
\(\Psi = \mathrm{e}^{(7/2)H\lvert\sigma\rvert}\chi(\sigma)\mathrm{e}^{-\mathrm{i}\omega x_4}\))

\[
\begin{aligned}
&\chi(\sigma) = \mathrm{e}^{-q_{\mathrm{b}}\lvert\sigma\rvert}\chi(0^+)\ \ (0 < \sigma \ll \sigma_{\mathrm{t}}),\qquad
\chi(\sigma) = \mathrm{e}^{-q_{\mathrm{b}}\lvert\sigma\rvert}\chi(0^-)\ \ (0 < -\sigma \ll \sigma_{\mathrm{t}}),\\
&\chi(0^+) = \exp\Big(\gamma^{(0)}\big(\mu_{\mathrm{L}}P_{\mathrm{L}} + \mu_{\mathrm{R}}P_{\mathrm{R}}\big)\Big)\chi(0^-),\qquad
\omega_0 = \omega_{\mathrm{r}} - \frac{\mathrm{i}}{2}\,\Gamma_\chi,\qquad \omega_{\mathrm{r}} \to \omega_{\mathrm{b}} = \sqrt{m_\Psi^2 - q_{\mathrm{b}}^2}
\end{aligned} \tag{FD-9}
\] \textbf{Symbols in (FD-9).}

\begin{itemize}
\tightlist
\item
  \(\chi(\sigma)\): The profile of a mode of dirac16complex along the
  hidden direction,
  \(\Psi = \mathrm{e}^{\frac72H\lvert\sigma\rvert}\chi(\sigma)\mathrm{e}^{\mathrm{i}(\vec k\cdot\vec x - \omega x_4)}\)
  (a 16-component spinor function of \(\sigma\)); as a four-dimensional
  field, \(\chi\) denotes the particle of the level of dirac16complex at
  the fold (a long-lived resonance) and \(\bar\chi\) its antiparticle.
  Unit: GeV\(^{2}\) (with the \(V_*\) normalisation). Status: DERIVED.
\item
  \(\sigma\): The proper coordinate of the hidden direction:
  \(\mathrm{d}\sigma = \lvert\cot z\rvert\,\mathrm{d}x_0\),
  \(\sigma = \mathrm{sgn}(\pi/2 - z)\ln(\sin z)/(6H)\) on one period
  \(0 < z < \pi\); \(\sigma = 0\) is the fold \(z = \pi/2\),
  \(\sigma \to \mp\infty\) the tips \(z \to 0, \pi\); one period of
  \(x_0\) is the whole line \(-\infty < \sigma < \infty\). Unit: length.
  Status: DERIVED.
\item
  \(\mathrm{e}\): Euler's number, the base of the natural logarithm;
  \(\mathrm{e}^{X}\) is the exponential of \(X\). Unit: dimensionless.
  Value: 2.718281828. Status: STANDARD.
\item
  \(q_{\mathrm{b}}\): The inverse width of a state of dirac16complex at
  the fold: inside the turning point its profile falls as
  \(\mathrm{e}^{-q_{\mathrm{b}}\lvert\sigma\rvert}\). Unit: GeV. Status:
  DERIVED.
\item
  \(\sigma_{\mathrm{t}}\): The turning point of a state of
  dirac16complex of the frequency \(\omega\): the depth at which the
  warped frequency equals the bulk mass,
  \(\lvert\omega\rvert\mathrm{e}^{H\sigma_{\mathrm{t}}} = m_\Psi\);
  beyond it every solution oscillates. Unit: length. Status: DERIVED.
\item
  \(\gamma^M,\ \gamma^{(a)}\): The Dirac matrices of dirac16complex
  along the coordinates, \(\gamma^M = \mathbf{e}_a{}^MT^a\),
  \(\gamma_M = g_{MN}\gamma^N\); \(\gamma^{(a)} = T^a\) along the frame
  direction \((a)\) (\(\gamma^{(0)}\) along \(\sigma\), \(\gamma^{(4)}\)
  along the time). Unit: dimensionless. Status: AUTHOR.
\item
  \(\mu_X,\ \mu_{\mathrm{L}},\ \mu_{\mathrm{R}}\): The dimensionless
  strengths of the fold for the left-handed and right-handed halves of
  dirac16complex, set by the mean fields of the Standard Model at the
  fold: the jump of the profile across the fold is the exponential of
  \(\gamma^{(0)}(\mu_{\mathrm{L}}P_{\mathrm{L}} + \mu_{\mathrm{R}}P_{\mathrm{R}})\);
  negative values attract. Unit: dimensionless. Status: DERIVED.
\item
  \(P_{\mathrm{L}}\): The author's left-handed chiral projector,
  \(P_{\mathrm{L}} = (\mathbf{1} - \Gamma_8)/2\). Unit: dimensionless.
  Status: AUTHOR.
\item
  \(P_{\mathrm{R}}\): The author's right-handed chiral projector,
  \(P_{\mathrm{R}} = (\mathbf{1} + \Gamma_8)/2\). Unit: dimensionless.
  Status: AUTHOR.
\item
  \(\omega_0\): The energy at rest (zero momentum along ordinary space
  and along the extra times) of a state of dirac16complex at the fold;
  complex for a resonance,
  \(\omega_0 = \omega_{\mathrm{r}} - \frac{\mathrm{i}}{2}\Gamma_\chi\).
  Unit: GeV. Status: DERIVED.
\item
  \(\omega_{\mathrm{r}}\): The mass of the particle \(\chi\): the real
  part of the energy at rest of the resonance of dirac16complex at the
  fold,
  \(\omega_{\mathrm{r}} = \omega_{\mathrm{b}}(1 - H/(2q_{\mathrm{b}})) + \dots\).
  Unit: GeV. Status: DERIVED.
\item
  \(\mathrm{i}\): The imaginary unit, \(\mathrm{i}^2 = -1\). Unit:
  dimensionless. Status: STANDARD.
\item
  \(\Gamma_\chi\): The width of the particle \(\chi\) at rest: the rate
  of its decay into the waves of dirac16complex in the bulk. Unit: GeV.
  Value: below \(1/t_0 = 1.5\times10^{-42}\) GeV for dark matter.
  Status: DERIVED.
\item
  \(\omega_{\mathrm{b}}\): The energy at zero momentum of the state of
  dirac16complex at the fold in the limit \(H \to 0\) (the bound state
  of UM-3): the particle \(\chi\) of UM-4 has the mass
  \(\omega_{\mathrm{r}} = \omega_{\mathrm{b}}(1 - H/(2q_{\mathrm{b}})) + \dots\),
  equal to \(\omega_{\mathrm{b}}\) up to the relative correction
  \(H/(2q_{\mathrm{b}})\). Unit: GeV. Status: DERIVED.
\item
  \(m_\Psi\): The mass parameter of dirac16complex in the bulk. Unit:
  GeV. Status: RECORD.
\item
  \(\Gamma_8\): The chirality matrix of eight dimensions,
  \(\Gamma_8 = T_0T_1\cdots T_7\) (the author's T16\^{}A{[}8{]})
  \(= \mathrm{diag}(-\mathbf{1}_8, \mathbf{1}_8)\);
  \(\Gamma_8^2 = \mathbf{1}\); it anticommutes with every \(T_a\). Unit:
  dimensionless. Status: AUTHOR.
\item
  \(\Psi,\ \bar\Psi\): dirac16complex: the author's 16-component complex
  spinor field of eight dimensions;
  \(\bar\Psi = \Psi^\dagger\sigma_{16}\) is its Dirac adjoint; in a term
  of the fold, \(\Psi\) is its value at \(\sigma = 0\). Unit:
  GeV\(^{7/2}\). Status: AUTHOR.
\item
  \(\mathbf{1},\ \mathbf{1}_8\): The identity matrix of the size fixed
  by the context (\(16 \times 16\) on dirac16complex, \(4 \times 4\) on
  a four-dimensional spinor, \(3 \times 3\) in generation or flavour
  space); \(\mathbf{1}_8\) is the \(8 \times 8\) identity matrix. Unit:
  dimensionless. Status: STANDARD.
\item
  \(\mathbf{e}_a{}^M,\ \mathbf{e}^a{}_M\): The vielbein (frame field) of
  the metric: \(g_{MN} = \eta_{ab}\,\mathbf{e}^a{}_M\mathbf{e}^b{}_N\);
  \(\mathbf{e}_a{}^M\) is its inverse;
  \(\mathbf{e}_{bN} = \eta_{bc}\mathbf{e}^c{}_N\); \(\mathbf{e}^{(a)}\)
  is the frame one-form \(\mathbf{e}^a{}_M\,\mathrm{d}X^M\). Unit:
  dimensionless. Status: STANDARD.
\item
  \(\omega\): A frequency: the energy of a mode
  \(\propto \mathrm{e}^{-\mathrm{i}\omega x_4}\) (\(\omega_o\) of the
  orbital \(o\)). Unit: GeV. Status: STANDARD.
\item
  \(\vec k\): The (comoving) momentum along ordinary space, \(\vec k^2\)
  its square. Unit: GeV. Status: STANDARD.
\item
  \(\vec x\): The position in ordinary space,
  \(\vec x = (x_1, x_2, x_3)\). Unit: length. Status: AUTHOR.
\item
  \(t_0\): The age of the universe. Unit: time. Value: 13.8 Gyr. Status:
  STANDARD.
\item
  \(\sigma_{16}\): The matrix of the Dirac adjoint,
  \(\sigma_{16} = T_0T_1T_2T_3\), real and symmetric; \(\sigma_{16}T_a\)
  is real and antisymmetric. Unit: dimensionless. Status: AUTHOR.
\item
  \({}^\dagger\): Hermitian conjugation: the complex conjugate of the
  transpose of a matrix or spinor; for an operator on the space of
  states, its adjoint with respect to the scalar product of that space
  (for the quantised dirac16complex the indefinite Krein scalar
  product). Status: STANDARD.
\item
  \(o\): The label of a single-particle state of dirac16complex: a
  Kohn-Sham orbital, or a mode of the quantised field; a sum over \(o\)
  runs over the orbitals or the modes. Status: STANDARD.
\end{itemize}

where \(\mu_{\mathrm{L}}\), \(\mu_{\mathrm{R}}\) are the strengths of
the fold for the left- and right-handed halves of dirac16complex (set by
the fields of the fold),
\(q_{\mathrm{b}} = -m_\Psi(\mu_{\mathrm{L}} + \mu_{\mathrm{R}})\tanh\varsigma/(2\varsigma)\),
\(\varsigma = \sqrt{\mu_{\mathrm{L}}\mu_{\mathrm{R}}}\)
(\(q_{\mathrm{b}} = -m_\Psi\mu_{\mathrm{L}}/2\) for
\(\mu_{\mathrm{R}} = 0\)), and \(\omega_{\mathrm{r}}\) tends to the
energy \(\omega_{\mathrm{b}}\) of the bound state of UM-3 when
\(H/q_{\mathrm{b}} \to 0\). Because the time is warped, the energy
\(\omega\) of the field enters its equation multiplied by
\(\mathrm{e}^{H\lvert\sigma\rvert}\): beyond the turning point
\(\sigma_{\mathrm{t}}\)
(\(\omega_{\mathrm{r}}\mathrm{e}^{H\sigma_{\mathrm{t}}} = m_\Psi\)) the
profile oscillates as a wave that leaves the fold, so the state is a
resonance of the width \(\Gamma_\chi\), which falls like
\(\mathrm{e}^{-2F(\upsilon_{\mathrm{b}})m_\Psi/H}\) with
\(F(\upsilon) = \ln((1 + \sqrt{1 - \upsilon^2})/\upsilon) - \sqrt{1 - \upsilon^2}\)
and \(\upsilon_{\mathrm{b}} = \omega_{\mathrm{b}}/m_\Psi\) (the
tunnelling through the region between the fold and
\(\sigma_{\mathrm{t}}\)). The profile is not a function of
\(\lvert\sigma\rvert\): it jumps at the fold,
\(\chi(0^+) \ne \chi(0^-)\). 8. \emph{The fold as a mirror for
dirac16complex.} The reflection across the fold acts on the profiles
\(\chi(\sigma)\) in two ways,

\[
(\mathsf{R}_1\chi)(\sigma) = T_0T_1T_2T_3\,\chi(-\sigma),\qquad (\mathsf{R}_2\chi)(\sigma) = T_0T_5T_6T_7\,\chi(-\sigma) \tag{FD-10}
\] \textbf{Symbols in (FD-10).}

\begin{itemize}
\tightlist
\item
  \(\mathsf{R}_1\): The reflection across the fold acting on a profile
  of dirac16complex together with the reflection of ordinary space,
  \((\mathsf{R}_1\chi)(\sigma) = T_0T_1T_2T_3\,\chi(-\sigma)\) (FD-10):
  a symmetry of dirac16complex, of gravity and of the mean field of the
  gas of the particles \(\chi\); the states at the fold of Krein charge
  \(-1\) are even, those of Krein charge \(+1\) odd. Unit:
  dimensionless. Status: DERIVED.
\item
  \(\mathsf{R}_2\): The reflection across the fold acting on a profile
  of dirac16complex together with the reflection of the extra times,
  \((\mathsf{R}_2\chi)(\sigma) = T_0T_5T_6T_7\,\chi(-\sigma)\) (FD-10):
  it maps the states at the fold of one Krein charge onto those of the
  other. Unit: dimensionless. Status: DERIVED.
\item
  \(\chi(\sigma)\): The profile of a mode of dirac16complex along the
  hidden direction,
  \(\Psi = \mathrm{e}^{\frac72H\lvert\sigma\rvert}\chi(\sigma)\mathrm{e}^{\mathrm{i}(\vec k\cdot\vec x - \omega x_4)}\)
  (a 16-component spinor function of \(\sigma\)); as a four-dimensional
  field, \(\chi\) denotes the particle of the level of dirac16complex at
  the fold (a long-lived resonance) and \(\bar\chi\) its antiparticle.
  Unit: GeV\(^{2}\) (with the \(V_*\) normalisation). Status: DERIVED.
\item
  \(\sigma\): The proper coordinate of the hidden direction:
  \(\mathrm{d}\sigma = \lvert\cot z\rvert\,\mathrm{d}x_0\),
  \(\sigma = \mathrm{sgn}(\pi/2 - z)\ln(\sin z)/(6H)\) on one period
  \(0 < z < \pi\); \(\sigma = 0\) is the fold \(z = \pi/2\),
  \(\sigma \to \mp\infty\) the tips \(z \to 0, \pi\); one period of
  \(x_0\) is the whole line \(-\infty < \sigma < \infty\). Unit: length.
  Status: DERIVED.
\item
  \(T_a,\ T^a\): The author's eight real \(16 \times 16\) Dirac matrices
  T16\^{}A{[}a{]} of 4+4.nb (\(a = 0, \dots, 7\); \(T_0, \dots, T_7\)),
  with \(\{T_a, T_b\} = 2\eta_{ab}\mathbf{1}\); \(T^a = \eta^{ab}T_b\).
  Unit: dimensionless. Status: AUTHOR.
\end{itemize}

(\(T_a\) the author's Dirac matrices; for states at rest
\(\mathsf{R}_1\) is the reflection across the fold combined with the
reflection of ordinary space, \(\mathsf{R}_2\) the reflection across the
fold combined with the reflection of the extra times). The Krein charge
of a state is its share of the conserved charge of dirac16complex, which
is indefinite: a state can carry the charge \(-1\) or \(+1\).
\(\mathsf{R}_1\) is a symmetry of dirac16complex, of gravity and of a
mean field of a gas of the particles \(\chi\) that is the same on both
sides of the fold: the four resonances of Krein charge \(-1\) (the
particles \(\chi\)) are even under it and the four of Krein charge
\(+1\) are odd, so such a mean field cannot mix the two signs of the
indefinite charge. \(\mathsf{R}_2\) maps the states of one Krein charge
onto those of the other and makes the two sets degenerate without the
self-interaction. 9. \emph{The planes of Horava and Witten.} In M-theory
(the eleven-dimensional theory whose limits include the superstring
theories) on an interval (a circle whose points are identified with
their mirror images), the two ten-dimensional planes left in place by
the reflection are mirror surfaces (property 1); each carries one
\(E_8\) gauge group while gravity fills the bulk. As in UM-4, gravity
fills the bulk; unlike UM-4, the gauge fields live on the mirror planes
(in UM-4 they live in the bulk), and the two sides of each plane are
identified (one physical side).

\textbf{What is not a fold.}

\begin{itemize}
\tightlist
\item
  \emph{The tips} \(z = 0\) and \(z = \pi\) (mod \(\pi\)): the mirror
  exchanges them within a period, but there the warp factor vanishes,
  \(g_{x_0x_0}\) diverges, and they lie at an infinite proper distance;
  they carry no thin source.
\item
  \emph{A smooth maximum of a warp factor}, such as
  \(1/\cosh^2(H\sigma)\): mirror symmetric and largest at
  \(\sigma = 0\), but without a kink and without a thin source.
\end{itemize}

\hypertarget{how-to-read-this-document}{%
\section{How to read this document}\label{how-to-read-this-document}}

This document records the origin of everything in the universe model
UM-4: the author's inputs and orders, the decisions taken while building
the model, the computations that verify every derived statement (each a
Python script that checks itself), the external measurements and
software that were used, and the commands that reproduce every result
and every document. It is a record of origins, decisions and
computations, not of the model itself, and so it contains no equations
other than those of the definition of the fold at its beginning.

The status labels used throughout UM-4: \textbf{AUTHOR} (from the
author), \textbf{RECORD} (from the record of this repository),
\textbf{STANDARD} (textbook physics or a published measurement),
\textbf{ASSUMED} (a choice without derivation), \textbf{INVENTED} (a new
interaction: the portals of dirac16complex), \textbf{DERIVED} (computed
and verified), \textbf{ESTIMATE} (an order of magnitude), \textbf{OPEN}
(not settled).

UM-4 is the successor of the universe model UM-3. The author's order of
2026-10-25, 13:15 PDT, which made UM-4, says ``keep UM-3 and create a
new UM-4, placed in new folders'': UM-3 is kept unchanged in the folder
\texttt{universe\_model\_3} of the repository, as committed in
\texttt{c5a93c4}, and everything of UM-4 is new, in the folder
\texttt{universe\_model\_4}, with programs of its own
(\texttt{um4\_*.py}) and documents of its own.

\hypertarget{the-authors-inputs}{%
\section{The author's inputs}\label{the-authors-inputs}}

\begin{longtable}[]{@{}
  >{\raggedright\arraybackslash}p{(\columnwidth - 2\tabcolsep) * \real{0.2683}}
  >{\raggedright\arraybackslash}p{(\columnwidth - 2\tabcolsep) * \real{0.7317}}@{}}
\toprule\noalign{}
\begin{minipage}[b]{\linewidth}\raggedright
input
\end{minipage} & \begin{minipage}[b]{\linewidth}\raggedright
content and place
\end{minipage} \\
\midrule\noalign{}
\endhead
\bottomrule\noalign{}
\endlastfoot
the order of 2026-10-24, 17:02 PDT & stop every earlier stage
permanently; a new revision folder; read \texttt{4+4.nb} and verify its
eight real \(16 \times 16\) Dirac matrices; the metric with \(b_4^2\)
(\(b_4\) a function of the time \(x_4\) in the extra-time components of
the metric, never identically 1); the complete coupled field equations
of the listed fields (Einstein-Lovelock gravity, the Higgs doublet with
the ad hoc vacuum, three generations of lepton doublets, the
right-handed electron, muon and tau, \(W^1, W^2, W^3, B\), three
generations of quark doublets, the gluons, dirac16complex coupled to
gravity with a non-trivial spin connection and with its left- and
right-handed halves coupled to every Standard-Model field, only the
left-handed half to the Higgs field); new phenomenological interactions;
an accurate approximate Kohn-Sham solution; the CMB, the CNB and the
BAO; the \(x_0\) mass mechanism; dark matter and dark energy; every
equation and every symbol in every document, never a reference to
another file; Markdown, LaTeX and PDF; push as soon as verified. In the
author's private prompt file (not committed). \\
the notebook 4+4.nb & the author's notebook ``eight, real-valued, 16 by
16, Dirac Matrices'' (Wolfram 15.0) with the matrices
\texttt{T16\^{}A{[}0..8{]}}; kept unchanged in the folder
\texttt{universe\_model\_3/author\_inputs} (its sha256 is given below)
and read from there by the programs of UM-4 (not copied) \\
the notebook of the record &
\texttt{Dirac\_Einstein-Lovelock\_4+4\_Nash.nb} in the root of the
repository: the coordinates \(X = \{x_0, \dots, x_7\}\) (Section
DEFCHART: \(x_0\) the hidden direction, \(x_1, x_2, x_3\) ordinary
space, \(x_4\) the time, \(x_5, x_6, x_7\) the extra times) and the
Einstein-Lovelock action with the weights \((1, 1/32, 1/1152)\) of the
Einstein, Gauss-Bonnet and cubic terms and the cosmological-constant
parameter \(\Lambda_{\mathrm{N}} = 10\) (In{[}134{]}); UM-4 uses this
action with the weights of the Gauss-Bonnet and cubic terms set to zero,
Einstein gravity with \(\Lambda = -21H^2\) \\
the decision of 2026-10-24, 17:32 & ``rather have three
\(\nu_{\mathrm{R}}\)'': three right-handed neutrinos, Dirac masses (in
the conversation) \\
the reason of 2026-10-24, 17:34 & Gogilashvili and Tamborra, ``Neutrino
Flavor Conversion Shapes the Rate of Failed Core-collapse Supernovae'',
arXiv:2605.16504: why the flavour conversion must be correct (a
third-party paper on the author's computer, read, not committed) \\
the standing rule of 2026-10-14 & no single calculation may take longer
than 423 s (7 min 3 s); every command runs under a time limit (recorded
in \texttt{claude\_failures.md} and \texttt{HANDOFF.md} of the
repository) \\
the orders of 2026-10-25, 03:21, 03:23, 03:54, 05:09 and 07:55 PDT &
given for UM-3 and carried over to UM-4: the sources on the fold (a thin
fluid and the gravity of the fold); the answers to the open questions
(the Krein norm, fields that depend on the extra times, particles faster
than light, the primordial era, the lightest neutrino mass and the Dirac
phase); one definition of the fold, with its examples, in every file
that uses the word; proper attribution of the stabiliser of the bulk of
UM-3 to Claude (UM-4 has no stabiliser); ``Gravitons must travel at c
\ldots{} my choices are wrong, and you MUST correct them'' \\
the order of 2026-10-25, 13:15 PDT & ``Since in the UM-3 model anything
that propagates in the bulk outruns light on the fold, and bulk scalars
and photons faster than light, my model is physically incorrect. You
must find a new metric model in which x4 is warped like the other
directions, both photons and gravitons actually propagate at c in the
bulk, photon waves are transverse spin-1 phenomena, Gravitational waves
(ripples in spacetime) are spin-2 phenomena (because they originate from
the rank-2 stress-energy tensor) and possess the following Polarization
Modes: the spin-2 nature of gravitational waves manifests physically as
characteristic''plus'' (+) and ``cross'' (×) polarization modes that
stretch and squeeze space perpendicularly. Find the best model and
re-calculate and re-document; however keep UM-3 and create a new UM-4,
placed in new folders.'' \\
\end{longtable}

The sha256 of the author's notebook \texttt{4+4.nb} as received and
committed:

\begin{Shaded}
\begin{Highlighting}[]
\NormalTok{7ec1709a0028bce0b443f116cf27b5d3e670de7aa195659a9c3274527270429d}
\end{Highlighting}
\end{Shaded}

The author's private files (his prompt files, his notes, the third-party
paper, his scratch notebook of 2026-10-16 with its \texttt{.mx} files
and the files \texttt{Generalized\_Kronecker\_Delta.*}) are never
committed; the repository's \texttt{.gitignore} excludes them.

\hypertarget{the-decisions-taken-in-um-4}{%
\section{The decisions taken in
UM-4}\label{the-decisions-taken-in-um-4}}

\begin{longtable}[]{@{}
  >{\raggedright\arraybackslash}p{(\columnwidth - 4\tabcolsep) * \real{0.3333}}
  >{\raggedright\arraybackslash}p{(\columnwidth - 4\tabcolsep) * \real{0.3333}}
  >{\raggedright\arraybackslash}p{(\columnwidth - 4\tabcolsep) * \real{0.3333}}@{}}
\toprule\noalign{}
\begin{minipage}[b]{\linewidth}\raggedright
decision
\end{minipage} & \begin{minipage}[b]{\linewidth}\raggedright
reason
\end{minipage} & \begin{minipage}[b]{\linewidth}\raggedright
status
\end{minipage} \\
\midrule\noalign{}
\endhead
\bottomrule\noalign{}
\endlastfoot
the time \(x_4\) warped by the factor \((\sin^2z)^{1/6}\) of the other
six directions of the fold, everything else of the author's metric kept
& the order of 13:15; the smallest change of the author's metric that
meets it: with the same factor the light cone of every slice of the bulk
is that of the fold, and among the static warps of the author's form
(ordinary space, the time and the extra times warped at their own rates)
only equal rates give a vacuum of gravity with a cosmological constant
(\texttt{um4\_geometry.py}, G8 and G10) & AUTHOR (the order), ASSUMED
(the way it is met), DERIVED (its properties) \\
Einstein gravity with \(\Lambda = -21H^2\) & the only choice of the
author's Lovelock weights for which gravitons travel on the light cone
on every background, not only around the static block
(\texttt{um4\_gravity.py}, L1 to L4); with it the static block is a
space of constant curvature \(-H^2\), a vacuum without a stabiliser, and
the fold a pure tension & DERIVED \\
the gauge fields of the Standard Model in the bulk & the order puts the
photon in the bulk, and with it the electroweak fields of which it is a
combination; the gluons are put there as well & AUTHOR (the photon),
ASSUMED (all gauge fields) \\
the Higgs doublet, the quarks and the leptons on the fold,
dirac16complex in the bulk & the author's specification of
dirac16complex as a field of eight dimensions; the fields of the
Standard Model other than the gauge fields are not described in the bulk
& ASSUMED \\
right-handed quark singlets and Yukawa couplings & masses of the quarks
and cancellation of the anomalies (defaults of the notification of
2026-10-24, 17:20; the author answered only on the neutrinos) &
ASSUMED \\
three right-handed neutrinos & the author's decision of 2026-10-24,
17:32 & AUTHOR \\
the fold fluid & the kink of the warp at the fold needs a thin source: a
pure tension \(\kappa\lambda_{\mathrm{F}} = 12H\) at the static block,
and a pressure \(q_{\mathrm{hold}}\) along the extra times that holds
their volume fixed (a free volume would change Newton's constant 8500
times faster than lunar laser ranging allows); its value today depends
on the pressure of the baryons along the extra times, taken as zero &
DERIVED (its values), ASSUMED (the medium, its holding pressure and the
pressure of the baryons), OPEN (the origin of the holding pressure, the
tuning of the tension) \\
the bound \(H > 2534\) GeV used in every document & the invisible width
of the Z boson, whose decay into the bulk the gauge fields of the bulk
allow (\texttt{um4\_propagators.py}, P5); the other bounds (contact
interactions, the parameter rho, Newton's law) are weaker & ESTIMATE \\
the portals of dirac16complex & the author's request for couplings of
both halves to every Standard-Model field, only the left-handed half to
the Higgs field & INVENTED \\
the pure-gauge \(W^2\) of the ad hoc vacuum & the Gauss law of
\(SU(2)\); in UM-4 it is the value at the fold of the field of the bulk
& DERIVED \\
no dependence of any field on the extra times & such fields are
tachyonic; the restriction is consistent but unstable, and a gauge
principle that removes the momentum along the extra times is needed &
ASSUMED (OPEN: the gauge principle) \\
the parameters of the Kohn-Sham computation (\(m_\Psi/H = 10^4\),
\(\mu_{\mathrm{L}} = -0.6\), \(\mu_{\mathrm{R}} = 0\), the box
\(\lvert\sigma\rvert \le 40/m_\Psi\) with the step \(0.03/m_\Psi\); the
weakly bound point \(q_{\mathrm{b}} = 0.03\,m_\Psi\),
\(H = 10^{-7}m_\Psi\)) & the states of dirac16complex at the fold are
resonances, so the gas is computed deep in the regime of long-lived
particles, far inside the turning point & ASSUMED \\
the lightest neutrino mass zero, the Dirac phase zero & not fixed by the
specification; the phase drops out of every result of UM-4, and DESI DR2
with the CMB and NuFIT 6.0 bounds the lightest mass below 23 meV &
ASSUMED (the mass within the bound), DERIVED (the phase is
irrelevant) \\
the primordial spectrum, the optical depth and the densities from Planck
2018 & the author's vacuum is not a solution; the primordial era of UM-4
can be an inflation of the fold with frozen extra times, which needs a
Weyl tensor of the bulk and a rate below \(H/10\); UM-4 does not specify
its inflaton & ASSUMED (inputs; the inflaton OPEN) \\
the abundance of the dark matter taken from the measurement & the
particles \(\chi\) of UM-4 are too heavy to be a thermal relic & ASSUMED
(the origin OPEN) \\
the quantisation of dirac16complex with positive energies & it makes the
vacuum stable; the modes of positive Krein charge are then ghosts, which
the chiral couplings of the fold produce; the classical stability of the
gas is not decided & DERIVED (the consequences), OPEN (the stability of
the gas, a unitary coupling) \\
one definition of the fold in every file & the order of 2026-10-25,
03:54: the definition with its examples,
\texttt{docs\_src/fold\_definition.md} for the documents and
\texttt{docs\_src/fold\_definition.txt} for the code and the JSON files,
inserted and checked by \texttt{tools/fold\_block.py}; the repository
files \texttt{README.md}, \texttt{HANDOFF.md} and
\texttt{claude\_failures.md} carry the definition of UM-4 as plain text
next to that of UM-3 & AUTHOR (the order) \\
the Kohn-Sham module in parts & the standing rule of 423 s: the module
runs in four parts with checkpoints, each under the time limit, and an
assembly that makes the checks & AUTHOR (the rule) \\
the lunar-laser-ranging bound of L. Biskupek, J. Müller and J.-M. Torre
(2021) & the bound on the change of Newton's constant used in every
document of UM-4 (two standard deviations) & STANDARD \\
\end{longtable}

\hypertarget{the-computations}{%
\section{The computations}\label{the-computations}}

Every derived statement rests on one of these scripts in
\texttt{universe\_model\_4/code/}. Each computes, checks itself and
writes its results to \texttt{universe\_model\_4/outputs/} and its
checks to \texttt{universe\_model\_4/reports/}.

\begin{longtable}[]{@{}
  >{\raggedright\arraybackslash}p{(\columnwidth - 4\tabcolsep) * \real{0.3333}}
  >{\raggedright\arraybackslash}p{(\columnwidth - 4\tabcolsep) * \real{0.3333}}
  >{\raggedright\arraybackslash}p{(\columnwidth - 4\tabcolsep) * \real{0.3333}}@{}}
\toprule\noalign{}
\begin{minipage}[b]{\linewidth}\raggedright
script
\end{minipage} & \begin{minipage}[b]{\linewidth}\raggedright
what it computes
\end{minipage} & \begin{minipage}[b]{\linewidth}\raggedright
checks
\end{minipage} \\
\midrule\noalign{}
\endhead
\bottomrule\noalign{}
\endlastfoot
\texttt{um4\_common.py} & the report of the checks (a library) & \\
\texttt{um4\_gammas.py} & the author's Dirac matrices, rebuilt from
\texttt{4+4.nb} (a library) & \\
\texttt{um4\_lovelock.py} & the Lovelock tensors from the frame
curvature (a library) & \\
\texttt{um4\_dirac\_matrices.py} & reality, Clifford relations,
Spin(4,4) and Pin(4,4), chiral projectors, agreement with every stored
output of the notebook and with the earlier records of the repository &
55 \\
\texttt{um4\_dirac\_matrices\_doc.py} & the document
\texttt{dirac\ matrices} from the checked results & \\
\texttt{um4\_geometry.py} & the proper chart with the warped time, the
exact curvature of the metric of UM-4, the static block of constant
curvature, its Lovelock tensors, the exact bulk equations of Einstein
gravity, the Ricci-flat folds, the pure tension, the light cone of every
slice, the spin connection with the author's matrices, the static warps
and their vacua & 10 \\
\texttt{um4\_gravity.py} & the vacua of the static family for any
weights, the graviton around the static block for any weights, the light
cone of the gravitons on every warped background, the choice of Einstein
gravity & 4 \\
\texttt{um4\_bulk.py} & the modes of the scalar and of the gauge field
of the bulk, the photon as a transverse wave of spin 1, the graviton of
the bulk and the gravitational wave with its plus and cross
polarisations, the bulk fermion, the boost invariance of dirac16complex,
the speeds, the speeds of UM-3 recomputed, the continuum of massive
modes from zero mass and its spectral density at the fold & 10 \\
\texttt{um4\_fold.py} & the bulk metric of an expanding fold to the
first order, the fold as a seven-dimensional Einstein brane, the exact
equations of the fold, the Friedmann equations with frozen extra times,
the size of the corrections, the Weyl term that dust requires & 6 \\
\texttt{um4\_foldsources.py} & the tension, the gravitational wave of
the fold, the induced gravity, the tuning of the tension, the fold fluid
& 5 \\
\texttt{um4\_background.py} & the vacua of the author's form, the
expanding fold, the free volume excluded by lunar laser ranging, the
fold factor, the ad hoc vacuum as a gauge rotation and its numbers &
6 \\
\texttt{um4\_propagators.py} & the propagators of the bulk from one
point of the fold to another, their expansions, Newton's and Coulomb's
laws, the W and Z resonances, the bounds on \(H\) & 5 \\
\texttt{um4\_causality.py} & no shortcut through the static bulk, the
light cone of an expanding fold, the largest advance, the closed
time-like curves & 4 \\
\texttt{um4\_fields.py} & the bilinears of dirac16complex, the Krein
operator, the electroweak vacuum with the ad hoc component, the
anomalies & 16 \\
\texttt{um4\_spectra.py} & the resonances of dirac16complex at the fold
with the exact matrices, their widths, their dispersion, the lifetime
condition, their Krein signature and their mirror (reads the outputs of
\texttt{um4\_propagators.py}) & 8 \\
\texttt{um4\_kohnsham.py} & the self-consistent Kohn-Sham states of a
gas of the particles \(\chi\), a weakly bound point, how dilute the gas
of the universe is & 10 \\
\texttt{um4\_krein.py} & the Krein form and its quantisation, the
symmetries and mirrors of the states at the fold, the flat directions of
the gas and its truncated response, the partner level, the selection
rules at the fold, stability and unitarity (reads the snapshots of
\texttt{um4\_kohnsham.py}) & 7 \\
\texttt{um4\_extratimes.py} & fields with a momentum along the extra
times: the dispersion, the states of dirac16complex with such a
momentum, the growth with frozen extra times, compact extra times and
their harmonics, the truncation & 5 \\
\texttt{um4\_neutrino.py} & the flavour blindness of the portals, the
matter conversion, the crossing criterion, the equipartition, the cosmic
neutrino background (CNB) & 5 \\
\texttt{um4\_flavour.py} & the neutrino flavours in the Sun and in
supernovae & 4 \\
\texttt{um4\_numass.py} & the lightest neutrino mass and the Dirac
phase: the masses of both orderings, the bounds of DESI DR2 and KATRIN,
the phase in the results of UM-4, the effect of the lightest mass, the
record against NuFIT 6.0 & 5 \\
\texttt{um4\_thermal.py} & the degrees of freedom, the relic density,
the portals and direct detection, the right-handed neutrinos, the ad hoc
vacuum before nucleosynthesis, no thermal relic, the binding by the
Higgs portal and the lifetime (reads the outputs of
\texttt{um4\_spectra.py}) & 8 \\
\texttt{um4\_primordial.py} & the author's vacuum (not a solution), the
exact de Sitter fold, the inflation with frozen extra times, its scales,
Higgs inflation (an extension of the action), the fold factor & 6 \\
\texttt{um4\_cosmology.py} & the cosmic microwave background (CMB), the
CNB and the baryon acoustic oscillations (BAO) with the Boltzmann code
CAMB and an independent background, against Planck 2018 and DESI DR2
(the second data release of the Dark Energy Spectroscopic Instrument) &
6 \\
\end{longtable}

The 20 programs have 185 checks, all passed. The documents are written
in Markdown and built with
\texttt{universe\_model\_4/tools/build\_docs.py} (pandoc to LaTeX,
pdflatex to PDF, deterministic).
\texttt{universe\_model\_4/docs\_src/registry.py} holds the meaning,
unit, value and status of every symbol;
\texttt{universe\_model\_4/tools/make\_inventory.py} writes the
inventory of the tagged equations from \texttt{FIELD\_EQUATIONS.md}, and
\texttt{universe\_model\_4/tools/symbols\_tool.py} copies every tagged
equation into the other documents (\texttt{sync}), writes the symbol
list after every display and the list of further symbols at the end of
every document (\texttt{expand}) and checks that every symbol of every
display is defined, that every symbol used in the meaning of a listed
symbol is listed in the same document, that every tagged equation that a
document cites is displayed in it, and that every tagged equation is
written identically in every document (\texttt{check}).
\texttt{universe\_model\_4/tools/symbols\_close.py} adds to the lists of
a document the entries that the meanings of its listed entries need
(with \texttt{universe\_model\_4/tools/closure\_check.py}), and
\texttt{universe\_model\_4/tools/inline\_symbols.py} checks that the
symbols of the inline formulas of the prose are listed in the document
as well (adding the missing entries to the list of further symbols; a
symbol without an entry, such as a label value or a generic letter of an
example, is defined in words where it is used).
\texttt{universe\_model\_4/tools/fold\_block.py} inserts the one
definition of the fold into every file that uses the word
(\texttt{insert}) and checks that every such file carries it word for
word (\texttt{check}). \texttt{universe\_model\_4/tools/nb\_reader.py}
reads the cells of the author's notebooks without Mathematica.

\hypertarget{external-measurements-and-software}{%
\section{External measurements and
software}\label{external-measurements-and-software}}

\begin{longtable}[]{@{}
  >{\raggedright\arraybackslash}p{(\columnwidth - 2\tabcolsep) * \real{0.5000}}
  >{\raggedright\arraybackslash}p{(\columnwidth - 2\tabcolsep) * \real{0.5000}}@{}}
\toprule\noalign{}
\begin{minipage}[b]{\linewidth}\raggedright
source
\end{minipage} & \begin{minipage}[b]{\linewidth}\raggedright
used for
\end{minipage} \\
\midrule\noalign{}
\endhead
\bottomrule\noalign{}
\endlastfoot
Planck Collaboration, Planck 2018 results VI, A\&A 641, A6 (2020), Table
2 (TT,TE,EE+lowE+lensing) & the cosmological parameters and the derived
quantities of the CMB \\
Planck 2018 plik-lite likelihood files (Planck Legacy Archive), via
github.com/heatherprince/planck-lite-py (Prince and Dunkley 2019) & the
CMB spectra (downloaded at run time into
\texttt{universe\_model\_4/external/}, not redistributed) \\
DESI Collaboration, DESI DR2 Results II, arXiv:2503.14738, Table IV &
the BAO distances; the 2.3 and 3.1 standard-deviation statements of its
abstract \\
W. Elbers et al.~(DESI Collaboration), arXiv:2503.14744, section IV C &
DESI DR2 with the CMB combined with NuFIT 6.0 (95 per cent,
cosmological-constant model): the lightest mass below 0.023 eV (normal
ordering) and 0.024 eV (inverted), the sum below 0.101 eV and 0.133 eV,
the probability 0.91 of the normal ordering; the bound 0.0642 eV on the
sum of three equal masses \\
ACT Collaboration (Atacama Cosmology Telescope), DR6 constraints on
extended cosmological models, arXiv:2503.14454; T. Louis et al.~(ACT
Collaboration), arXiv:2503.14452 & the effective number of neutrino
species \(N_{\mathrm{eff}} = 2.89 \pm 0.11\), the sum of the neutrino
masses below 0.082 eV; the spectral index \(0.974 \pm 0.003\) (ACT DR6
with Planck, CMB lensing and DESI) \\
Planck Collaboration, Planck 2018 results X, A\&A 641, A10 (2020); M.
Tristram et al., Phys. Rev.~D 105, 083524 (2022) (Planck PR4 with
BICEP/Keck 2018 and BAO) & the spectral index \(0.9649 \pm 0.0042\); the
bound \(r_{\mathrm{T}} < 0.032\) (95 per cent) on the tensor-to-scalar
ratio \\
P. Janot and S. Jadach, Phys. Lett. B 803, 135319 (2020) & the number of
light neutrinos \(2.9963 \pm 0.0074\) measured at the Z boson, which
bounds the decay of the Z boson into the bulk (\(H > 2534\) GeV) \\
ATLAS Collaboration, JHEP 11 (2020) 005 & the limit
\(\Lambda_{\mathrm{CI}} > 26.0\) TeV on the scale of contact
interactions of electrons and muons (\(H > 1.3\) TeV) \\
Particle Data Group, Review of Particle Physics (2022 and 2024) & the
masses of the Higgs boson (125.25 GeV), of the quarks and charged
leptons and of the nucleon, the Fermi constant
\(1.1663788\times10^{-5}\) GeV\(^{-2}\) and \(v = 246.22\) GeV, the
fine-structure constant; the parameter rho of the electroweak fits,
\(1.00031 \pm 0.00019\) (2024) \\
J. G. Lee, E. G. Adelberger, T. S. Cook, S. M. Fleischer and B. R.
Heckel, Phys. Rev.~Lett. 124, 101101 (2020) & Newton's law tested by
torsion balances down to 52 micrometres \\
B. P. Abbott et al.~(LIGO Scientific, Virgo, Fermi GBM and INTEGRAL
Collaborations), Astrophys. J. Lett. 848, L13 (2017) & the speed of
gravitational waves,
\(-3\times10^{-15} < c_{\mathrm{g}} - 1 < 7\times10^{-16}\) (GW170817
and GRB 170817A) \\
L. Biskupek, J. Müller and J.-M. Torre, Universe 7, 34 (2021) & the
bound \((-5.0 \pm 9.6)\times10^{-15}\) per year of lunar laser ranging
on the change of Newton's constant, against a free volume of the extra
times \\
W. Israel, Nuovo Cimento B 44, 1 (1966); T. Shiromizu, K. Maeda and M.
Sasaki, Phys. Rev.~D 62, 024012 (2000); L. Randall and R. Sundrum, Phys.
Rev.~Lett. 83, 4690 (1999) & the junction condition of the fold; the
exact equations of the fold in a bulk of constant curvature; the brane
in anti-de Sitter space whose tension and fine tuning the fold of UM-4
shares \\
A. D. Sakharov, Dokl. Akad. Nauk SSSR 177, 70 (1967); D. V. Vassilevich,
Phys. Rep.~388, 279 (2003); S. N. Solodukhin, arXiv:1502.03758 & induced
gravity; the heat-kernel coefficients of scalars, fermions and gauge
fields \\
K. Griest and M. Kamionkowski, Phys. Rev.~Lett. 64, 615 (1990) & the
unitarity bound on the mass of a thermal relic \\
LUX-ZEPLIN (LZ) Collaboration, Phys. Rev.~Lett. 135, 011802 (2025),
arXiv:2410.17036 version 3 (4.2 tonne-years) & the limit
\(2.2\times10^{-48}\) cm\(^2\) at 40 GeV (90 per cent confidence) on
spin-independent scattering \\
Borexino Collaboration, Nature 562, 505 (2018) & the survival
probabilities of solar neutrinos \\
the global fit NuFIT 5.2 of the neutrino oscillation data (Esteban,
Gonzalez-Garcia, Maltoni, Schwetz and Zhou, JHEP 09 (2020) 178; version
5.2 of November 2022 at www.nu-fit.org), normal ordering, the fit
without the Super-Kamiokande atmospheric data &
\(\sin^2\theta_{12} = 0.303\), \(\sin^2\theta_{13} = 0.02203\),
\(\sin^2\theta_{23} = 0.572\), \(\Delta m^2_{21} = 7.41\times10^{-5}\)
eV\(^2\), \(\Delta m^2_{31} = 2.511\times10^{-3}\) eV\(^2\) \\
the global fit NuFIT 6.0 (I. Esteban, M. C. Gonzalez-Garcia, M. Maltoni,
I. Martinez-Soler, J. P. Pinheiro and T. Schwetz, JHEP 12 (2024) 216,
arXiv:2410.05380), data set IC24 with the Super-Kamiokande atmospheric
data, and IC19 without it (\(\sin^2\theta_{12} = 0.307\),
\(\sin^2\theta_{23} = 0.561\), \(\sin^2\theta_{13} = 0.02195\),
\(\delta_{\mathrm{CP}} = 177^\circ\),
\(\Delta m^2_{31} = 2.534\times10^{-3}\) eV\(^2\)) &
\(\sin^2\theta_{12} = 0.308\), \(\sin^2\theta_{23} = 0.470\),
\(\sin^2\theta_{13} = 0.02215\), \(\delta_{\mathrm{CP}} = 212^\circ\),
\(\Delta m^2_{21} = 7.49\times10^{-5}\) eV\(^2\),
\(\Delta m^2_{31} = 2.513\times10^{-3}\) eV\(^2\) (normal ordering), and
the inverted ordering, for the comparison with the record \\
KATRIN Collaboration, Science 388, 180 (2025) & the limit 0.45 eV (90
per cent) on the effective electron-neutrino mass of tritium beta
decay \\
Long, Lunardini and Sabancilar, arXiv:1405.7654 & the capture rate of
the CNB on tritium (4 and 8 per year per 100 g for Dirac and Majorana
neutrinos) \\
Gogilashvili and Tamborra, arXiv:2605.16504 (version 3) & the flavour
equipartition (the equal sharing of the three flavours that fast flavour
conversion produces) and its effect on the explosions of massive
stars \\
Bahcall's exponential solar density profile; Petcov's crossing
probability; Izaguirre, Raffelt and Tamborra's dispersion relation;
Morinaga's crossing theorem; Ma and Bertschinger's Boltzmann equations;
F. Bezrukov and M. Shaposhnikov, Phys. Lett. B 659, 703 (2008) (Higgs
inflation) & the standard physics named in the documents \\
CAMB 2.0.5 (Lewis, Challinor and Lasenby 2000) & the Boltzmann hierarchy
and the CMB spectra \\
Python 3.13.17 with numpy 2.5.3, scipy 1.18.1, sympy 1.14.0, mpmath
1.3.0, matplotlib 3.11.2; pandoc 3.1.3; pdfTeX (TeX Live 2023) & every
computation and document \\
\end{longtable}

\hypertarget{how-to-reproduce}{%
\section{How to reproduce}\label{how-to-reproduce}}

From the root of a clone of the repository, with the software above
installed, each command run under \texttt{timeout\ 423}:

\begin{Shaded}
\begin{Highlighting}[]
\NormalTok{python universe\_model\_4/code/um4\_dirac\_matrices.py}
\NormalTok{python universe\_model\_4/code/um4\_dirac\_matrices\_doc.py}
\NormalTok{python universe\_model\_4/code/um4\_geometry.py}
\NormalTok{python universe\_model\_4/code/um4\_gravity.py}
\NormalTok{python universe\_model\_4/code/um4\_bulk.py}
\NormalTok{python universe\_model\_4/code/um4\_fold.py}
\NormalTok{python universe\_model\_4/code/um4\_foldsources.py}
\NormalTok{python universe\_model\_4/code/um4\_background.py}
\NormalTok{python universe\_model\_4/code/um4\_propagators.py}
\NormalTok{python universe\_model\_4/code/um4\_causality.py}
\NormalTok{python universe\_model\_4/code/um4\_fields.py}
\NormalTok{python universe\_model\_4/code/um4\_spectra.py}
\NormalTok{python universe\_model\_4/code/um4\_kohnsham.py part K1}
\NormalTok{python universe\_model\_4/code/um4\_kohnsham.py part I}
\NormalTok{python universe\_model\_4/code/um4\_kohnsham.py part K6}
\NormalTok{python universe\_model\_4/code/um4\_kohnsham.py part K9}
\NormalTok{python universe\_model\_4/code/um4\_kohnsham.py assemble}
\NormalTok{python universe\_model\_4/code/um4\_krein.py}
\NormalTok{python universe\_model\_4/code/um4\_extratimes.py}
\NormalTok{python universe\_model\_4/code/um4\_neutrino.py}
\NormalTok{python universe\_model\_4/code/um4\_flavour.py}
\NormalTok{python universe\_model\_4/code/um4\_numass.py}
\NormalTok{python universe\_model\_4/code/um4\_thermal.py}
\NormalTok{python universe\_model\_4/code/um4\_primordial.py}
\NormalTok{python universe\_model\_4/code/um4\_cosmology.py}
\NormalTok{python universe\_model\_4/docs\_src/registry.py}
\NormalTok{python universe\_model\_4/tools/fold\_block.py insert}
\NormalTok{python universe\_model\_4/tools/make\_inventory.py}
\NormalTok{python universe\_model\_4/tools/symbols\_tool.py sync}
\NormalTok{python universe\_model\_4/tools/symbols\_tool.py check}
\NormalTok{python universe\_model\_4/tools/fold\_block.py check}
\NormalTok{python universe\_model\_4/tools/build\_docs.py}
\end{Highlighting}
\end{Shaded}

Every script prints one line per check and ends with ``n/n checks
passed''. \texttt{um4\_spectra.py} reads the outputs of
\texttt{um4\_propagators.py}, \texttt{um4\_thermal.py} those of
\texttt{um4\_spectra.py}, \texttt{um4\_krein.py} the snapshots of the
parts of \texttt{um4\_kohnsham.py},
\texttt{um4\_dirac\_matrices\_doc.py} the outputs of
\texttt{um4\_dirac\_matrices.py}, and \texttt{um4\_cosmology.py}
downloads the Planck files once (about 7 MB). By the standing rule no
single calculation may take longer than 423 seconds: a part of
\texttt{um4\_kohnsham.py} that reaches its own limit of 400 seconds
saves its state and exits with the code 3, and is then simply run again
until it exits with 0. The documents are rebuilt identically (no dates
in the LaTeX, fixed PDF metadata).

\hypertarget{further-symbols}{%
\section{Further symbols}\label{further-symbols}}

\textbf{Further symbols.} These symbols appear in the definitions of the
symbols listed after the equations of this document.

\begin{itemize}
\tightlist
\item
  \(M_{\mathrm{red}}\): The reduced Planck mass,
  \(M_{\mathrm{red}} = (8\pi G_{\mathrm{N}})^{-1/2}\). Unit: GeV. Value:
  \(2.435\times10^{18}\) GeV. Status: STANDARD.
\item
  \(\nu_\alpha,\ \nu_i\): A neutrino field: \(\nu_\alpha\) the flavour
  states (the flavour index runs over the electron, muon and tau
  flavours), \(\nu_i\) the mass states (\(i = 1, 2, 3\)). Unit:
  GeV\(^{3/2}\). Status: STANDARD.
\item
  \(\rho\): The total energy density of the matter of the fold (all
  species); \(\Delta\rho\) its perturbation. Unit: GeV\(^4\). Status:
  STANDARD.
\item
  \(B_M,\ B_\mu,\ B_{MN}\): The hypercharge gauge field of \(U(1)_Y\):
  in the bulk \(B_M\) (geometric normalisation), on the fold the
  canonically normalised four-dimensional field \(B_\mu(x)\) of its
  massless mode, \(g_1B_\mu(x) = B_\mu(\sigma = 0, x)\); \(B_{MN}\),
  \(B_{\mu\nu}\) their field strengths. Unit: GeV (field), GeV\(^2\)
  (field strength). Status: AUTHOR.
\item
  \(W^I_M,\ W^I_\mu,\ W^I_{MN}\): The gauge field of \(SU(2)_L\)
  (\(I = 1, 2, 3\)): in the bulk \(W^I_M\) (geometric normalisation), on
  the fold the canonically normalised four-dimensional field
  \(W^I_\mu(x)\) of its massless mode,
  \(g_2W^I_\mu(x) = W^I_\mu(\sigma = 0, x)\); \(W^I_{MN}\),
  \(W^I_{\mu\nu}\) their field strengths;
  \(W^\pm_\mu = (W^1_\mu \mp \mathrm{i}W^2_\mu)/\sqrt2\). Unit: GeV
  (field), GeV\(^2\) (field strength). Status: AUTHOR.
\item
  \(m\): A mass: \(m_i\) the mass of the neutrino mass state \(i\);
  \(\Delta m_{ij}^2 = m_i^2 - m_j^2\) the squared-mass difference of the
  states \(i\) and \(j\) (\(\Delta m_{21}^2\), \(\Delta m_{31}^2\),
  \(\Delta m_{32}^2\)); \(\Delta m^2\) without labels the difference of
  the two states of a resonance, \(m_2^2 - m_1^2\), multiplied by
  \(v_{\mathrm{loc}}^2/v^2\) where this differs from 1; \(m\) without a
  label the mass of a mode of a field of the bulk as the fold sees it,
  \(m^2 = \omega^2 - k^2\); for the W and Z bosons their complex pole
  mass, \(m^2 = m_0^2 - \mathrm{i}m_0\Gamma_{\mathrm{bulk}}\) to the
  leading order; in the relic abundance of the particles \(\chi\) their
  mass. Unit: GeV. Status: STANDARD.
\item
  \(v\): The vacuum expectation value of the lower component of the
  Higgs doublet, \(v^2 = -\mu^2/\lambda\). Unit: GeV. Value: 246.22 GeV
  \(= (\sqrt2\,G_{\mathrm{F}})^{-1/2}\) with the Fermi constant
  \(G_{\mathrm{F}} = 1.1663788\times10^{-5}\) GeV\(^{-2}\) of the muon
  lifetime (Particle Data Group). Status: STANDARD.
\item
  \(\mu^2\): The mass parameter of the Higgs potential, \(\mu^2 < 0\)
  (it enters only as \(\mu^2\)). Unit: GeV\(^2\). Value:
  \(\mu^2 = -\lambda v^2 = -m_h^2/2 = -(88.6\ \mathrm{GeV})^2\). Status:
  AUTHOR.
\item
  \(k\): A wave number: the momentum of a plane wave along \(x_3\),
  \(\mathrm{e}^{\mathrm{i}(kx_3 - \omega x_4)}\); in the cosmological
  perturbations the comoving wave number of a Fourier mode,
  \(k = \lvert\vec k\rvert\). Unit: GeV (1/length). Status: STANDARD.
\item
  \(j\): A second generation or mass-state index, or a second direction
  of ordinary space (as \(i\)). Status: STANDARD.
\item
  \(G_{\mathrm{F}}\): The Fermi constant,
  \(G_{\mathrm{F}} = 1/(\sqrt2\,v^2)\) in the Standard-Model vacuum (its
  measured value); in the equations of the neutrino flavours it stands
  for the local Fermi constant of the ad hoc vacuum. Unit: GeV\(^{-2}\).
  Value: \(1.1664\times10^{-5}\) GeV\(^{-2}\). Status: STANDARD.
\item
  \(I\): A weak-isospin (adjoint SU(2)) index running over \(1, 2, 3\);
  a repeated index is summed. Status: STANDARD.
\item
  \(\Delta\): A change, difference or small variation of the quantity
  that follows (for example an interval of the time, a perturbation of a
  density, a difference of two squared masses, or an addition to an
  effective number of species). Status: STANDARD.
\item
  \(\Gamma_{\mathrm{bulk}}\): The width of the W or the Z boson into the
  bulk (its decay into the massive modes of the gauge fields of the
  bulk), \(\Gamma_{\mathrm{bulk}} = m_0(m_0/H)^3/3\). Unit: GeV. Value:
  1.4 MeV for the Z and 0.9 MeV for the W at \(H = 2.5\) TeV. Status:
  DERIVED.
\item
  \(\lambda\): The quartic coupling of the Higgs potential,
  \(\lambda > 0\). Unit: dimensionless. Value:
  \(\lambda = m_h^2/(2v^2) = 0.129\). Status: AUTHOR.
\item
  \(g_1\): The \(U(1)_Y\) gauge coupling of the four-dimensional world
  (usually \(g'\); not normalised for grand unification): the coupling
  of the massless mode of the hypercharge field of the bulk,
  \(g_1^2 = 3H\mathsf{g}_1^2/(2V_*)\). Unit: dimensionless. Value:
  \(g_1(m_Z) = 0.3574\). Status: STANDARD.
\item
  \(g_2\): The \(SU(2)_L\) gauge coupling of the four-dimensional world
  (usually \(g\)): the coupling of the massless mode of the \(SU(2)_L\)
  field of the bulk, \(g_2^2 = 3H\mathsf{g}_2^2/(2V_*)\). Unit:
  dimensionless. Value: \(g_2(m_Z) = 0.6517\). Status: STANDARD.
\item
  \(i\): A generation index (\(i = 1, 2, 3\)), which also labels the
  three neutrino mass states; in a frame component such as
  \(\mathcal{S}^{(i)}{}_{(i)}\) one of the three directions of ordinary
  space (the three components are equal), and in a tensor of ordinary
  space such as \(\mathsf{h}_{ij}\) one of the directions
  \(x_1, x_2, x_3\); a repeated index is not summed unless a sum sign is
  written. Status: STANDARD.
\item
  \(m_0\): The mass of the W or the Z boson without the massive modes of
  the bulk (the mass of the massless mode with the Higgs vacuum),
  \(m_0^2 = \frac32HM_{\mathrm{loc}}^2\);
  \(m_0^2 = g_2^2(v^2 + \mathcal{A}^2)/4\) for the W. Unit: GeV. Status:
  DERIVED.
\item
  \(m_u, \dots, m_3\): The masses of the quarks (\(m_u, m_c, m_t\);
  \(m_d, m_s, m_b\)), of the charged leptons (\(m_e, m_\mu, m_\tau\))
  and of the three neutrino mass states (\(m_1, m_2, m_3\), Dirac
  masses). Unit: GeV. Value: measured values;
  \(\Delta m_{21}^2 = 7.41\times10^{-5}\) eV\(^2\),
  \(\lvert\Delta m_{31}^2\rvert = 2.511\times10^{-3}\) eV\(^2\) in the
  computations of UM-4 (the global fit NuFIT 5.2 of the oscillation
  data, November 2022), \(7.49\times10^{-5}\) and \(2.513\times10^{-3}\)
  eV\(^2\) in NuFIT 6.0 (2024); the lightest mass \(m_1 = 0\) in the
  computations, below 23 meV by DESI DR2 with the CMB and NuFIT 6.0.
  Status: STANDARD.
\item
  \(m_h\): The mass of the Higgs boson. Unit: GeV. Value: 125.25 GeV
  (Particle Data Group 2022). Status: STANDARD.
\item
  \(v_{\mathrm{loc}}\): The local value of the electroweak vacuum,
  \(v_{\mathrm{loc}}^2 = 2\langle\Phi^\dagger\Phi\rangle\): it sets the
  local masses of the quarks and leptons (each mass times
  \(v_{\mathrm{loc}}/v\)) and the local Fermi constant;
  \(v_{\mathrm{loc}} = v\) in the Standard-Model vacuum. Unit: GeV.
  Value: \(v\) in the late universe outside dense dirac16complex.
  Status: DERIVED.
\item
  \(M_{\mathrm{loc}}\): The mass term that the Higgs vacuum on the fold
  gives a gauge field of the bulk:
  \(M_{\mathrm{loc}}^2\delta_{\mathrm{D}}(\sigma)\) in its profile
  equation; for the W,
  \(M_{\mathrm{loc}}^2 = \mathsf{g}_2^2(v^2 + \mathcal{A}^2)/(4V_*)\).
  Unit: GeV\(^{1/2}\) (\(M_{\mathrm{loc}}^2\) in GeV). Status: DERIVED.
\item
  \(\Phi\): The Higgs doublet (two complex components, hypercharge
  \(1/2\)), a field of the fold. Unit: GeV. Status: AUTHOR.
\item
  \(\mathcal{A}\): The author's ad hoc upper component of the Higgs
  vacuum, almightyAdHoc,
  \(\mathcal{A}(x_4) = v_{\mathcal{A}}\mathrm{e}^{-2a_{40}x_4}\). Unit:
  GeV. Status: AUTHOR.
\item
  \(\mathsf{g}_1,\ \mathsf{g}_2,\ \mathsf{g}_3\): The gauge couplings of
  eight dimensions of \(U(1)_Y\), \(SU(2)_L\) and \(SU(3)_C\): the
  coefficients of the kinetic terms of the gauge fields of the bulk in
  their geometric normalisation; the four-dimensional couplings are
  \(g_k^2 = 3H\mathsf{g}_k^2/(2V_*)\). Unit: GeV\(^{-2}\)
  (\(\mathsf{g}^2\): length\(^4\)). Status: ASSUMED.
\item
  \(\mathsf{h}_{ij}\): The transverse traceless perturbation of the
  metric of ordinary space by a gravitational wave as the
  four-dimensional world sees it:
  \(\hat h_{ij} = \mathrm{e}^{2a_4}(\delta_{ij} + \mathsf{h}_{ij})\)
  (\(i, j = 1, 2, 3\)). Unit: dimensionless. Status: STANDARD.
\item
  \(v_{\mathcal{A}}\): The amplitude of almightyAdHoc at \(x_4 = 0\);
  the author's almightyAdHoc, \(\mathrm{e}^{-2a_{40}x_4}\) (written
  Exp{[}-2 a40 x4{]} in his notebook), has \(v_{\mathcal{A}} = 1\) in
  the units of his notebook. Unit: GeV. Value: 1 in the units of the
  author's notebook; UM-4 keeps it free (the examples take
  \(v_{\mathcal{A}} = v\) and \(10v\)). Status: AUTHOR (the form),
  ASSUMED (a free amplitude).
\item
  \(N_{\mathrm{eff}}\): The effective number of neutrino species: the
  energy density of the relativistic particles other than photons in
  units of that of one neutrino species whose temperature is
  \((4/11)^{1/3}\) times that of the photons. Unit: dimensionless.
  Value: 3.044 (Standard Model); measured \(2.99 \pm 0.17\) (Planck
  2018). Status: STANDARD.
\item
  \(\mathcal{T}\): The energy-momentum tensor of any one term of the
  fold (in \(\mathcal{T}^r{}_r = \mathcal{L}_{\mathrm{fold}}\)). Unit:
  GeV\(^4\). Status: STANDARD.
\item
  \(\mathcal{W}\): The warp factor: in a warped metric, the function of
  the position along one direction that multiplies the squared steps of
  other directions (its square root multiplies their lengths); in UM-4
  the function of \(\sigma\) that multiplies all seven directions of the
  fold, ordinary space, the time and the extra times,
  \(\mathcal{W} = \mathrm{e}^{-2H\lvert\sigma\rvert} = (\sin^2z)^{1/6}\),
  equal to 1 on the fold. Unit: dimensionless. Status: DERIVED.
\item
  \(c_{\mathrm{g}}\): The speed of gravitational waves on the fold, in
  units of the speed of light. Unit: dimensionless. Value: 1 at every
  wavelength for the four-dimensional graviton on the static block and
  on every fold with frozen extra times, and for the massless graviton
  of the bulk at every depth; measured
  \(-3\times10^{-15} \le c_{\mathrm{g}} - 1 \le 7\times10^{-16}\)
  (GW170817). Status: DERIVED.
\item
  \(\delta_{\mathrm{CP}}\): The Dirac phase of the mixing matrix
  \(U_{\mathrm{PMNS}}\). Unit: degree. Value: \(212^\circ\) (NuFIT 6.0,
  normal ordering); 0 in the computations of UM-4. Status: STANDARD.
\item
  \(\kappa_7\): The gravitational coupling of the fold: to the first
  order in its curvature over \(H^2\) the fold obeys the Einstein
  equations of seven dimensions
  \(\hat G^a{}_b = \kappa_7\mathsf{t}^a{}_b\) with
  \(\kappa_7 = 5\kappa H/2\), the coupling of the massless graviton of
  the bulk bound to the fold; four-dimensional gravity on the fold has
  \(8\pi G_{\mathrm{N}} = \kappa_7/V_*\) (plus the induced term, which
  is negligible). Unit: GeV\(^{-5}\). Value: \(5\kappa H/2\). Status:
  DERIVED.
\item
  \(\nu_{\mathrm{R}}\): The three right-handed neutrinos (gauge
  singlets). Unit: GeV\(^{3/2}\). Status: ASSUMED.
\item
  \(q_{\mathrm{hold}}\): The pressure along the extra times that the
  fold fluid adds to its tension to hold the volume of the extra times
  fixed (four-dimensional normalisation). Unit: GeV\(^4\). Value:
  \(-(\rho - 3\mathcal{P})/2 - q_{\mathrm{SM}} - q_\chi - q_{\mathrm{P}}\):
  0 for radiation, \(-\rho/2\) for dust, \(-\rho\) for a vacuum energy;
  today \(-0.842\,\rho_{\mathrm{c}}^{(0)}\) if the baryons have no
  pressure along the extra times (ASSUMED). Status: DERIVED (the value),
  ASSUMED (that the fold fluid supplies it).
\item
  \(r_{\mathrm{T}}\): The ratio of the amplitudes of the primordial
  spectra of gravitational waves (tensor) and of curvature (scalar) at
  the pivot scale. Unit: dimensionless. Value: below 0.032 (Planck PR4
  with BICEP/Keck 2018 and BAO, M. Tristram et al.~2022, 95 per cent).
  Status: STANDARD.
\item
  \(\mathsf{t}^a{}_b,\ \mathsf{t}\): The matter of the fold above its
  tension: the thin source minus the tension,
  \(\mathcal{S}^a{}_b = -\lambda_{\mathrm{F}}\delta^a{}_b + \mathsf{t}^a{}_b\)
  (eight-dimensional normalisation); \(\mathsf{t} = \mathsf{t}^c{}_c\)
  its trace; with frozen extra times
  \(\mathsf{t}^a{}_b = \mathrm{diag}(P_7, P_7, P_7, -\rho_7, q_7, q_7, q_7)\).
  Unit: GeV\(^7\). Status: DERIVED.
\item
  \(\theta_{12},\ \theta_{13},\ \theta_{23}\): The three mixing angles
  of \(U_{\mathrm{PMNS}}\) (its fourth parameter is the Dirac phase,
  which does not enter the electron-flavour results of UM-4). Unit:
  dimensionless. Value: \(33.40^\circ\), \(8.54^\circ\), \(49.14^\circ\)
  in the computations of UM-4 (\(\sin^2\theta_{12} = 0.303\),
  \(\sin^2\theta_{13} = 0.02203\), \(\sin^2\theta_{23} = 0.572\): NuFIT
  5.2, November 2022, normal ordering, the fit without the
  Super-Kamiokande atmospheric data); \(33.71^\circ\), \(8.56^\circ\),
  \(43.28^\circ\) in NuFIT 6.0 (2024, with the Super-Kamiokande
  atmospheric data; \(\sin^2\theta_{12} = 0.308\),
  \(\sin^2\theta_{13} = 0.02215\), \(\sin^2\theta_{23} = 0.470\)).
  Status: STANDARD.
\item
  \(\upsilon,\ \upsilon_{\mathrm{b}}\): The argument of the tunnelling
  exponent \(F\), a ratio of an energy to the bulk mass \(m_\Psi\);
  \(\upsilon_{\mathrm{b}} = \omega_{\mathrm{b}}/m_\Psi\) that of the
  particle \(\chi\). Unit: dimensionless. Status: DERIVED.
\item
  \(F(\upsilon)\): The tunnelling exponent of the width of the particles
  \(\chi\) per unit of \(m_\Psi/H\): the integral of the local decay
  rate \(q\) over the region inside the turning point, divided by
  \(m_\Psi/H\), as a function of the ratio \(\upsilon\) of the energy to
  the bulk mass. Unit: dimensionless. Value: 0.049 for
  \(\upsilon = 0.866\); \((q_{\mathrm{b}}/\omega_{\mathrm{b}})^3/3\) for
  weak binding. Status: DERIVED.
\item
  \(P_7\): The pressure along ordinary space of the matter of the fold
  above its tension in the seven-dimensional normalisation (the
  component \(\mathsf{t}^{(i)}{}_{(i)}\)). Unit: GeV\(^7\). Status:
  DERIVED.
\item
  \(U_{\mathrm{PMNS}}\): The lepton mixing matrix
  (Pontecorvo-Maki-Nakagawa-Sakata) of the three Dirac neutrinos; its
  entries \((U_{\mathrm{PMNS}})_{\alpha i}\) connect the flavour state
  \(\nu_\alpha\) with the mass state \(\nu_i\). Unit: dimensionless.
  Value: \(\theta_{12} = 33.40^\circ\), \(\theta_{23} = 49.14^\circ\),
  \(\theta_{13} = 8.54^\circ\) (NuFIT 5.2, November 2022, normal
  ordering, the fit without the Super-Kamiokande atmospheric data); the
  Dirac phase is set to zero (ASSUMED; the fit prefers \(197^\circ\)).
  Status: STANDARD, ASSUMED (the phase).
\item
  \(\mathcal{L}_{\mathrm{fold}}\): The Lagrangian of any one term of the
  fold (four-dimensional normalisation). Unit: GeV\(^4\). Status:
  STANDARD.
\item
  \(\mathcal{P}\): The total pressure of the matter of the fold;
  \(\Delta\mathcal{P}\) its perturbation. Unit: GeV\(^4\). Status:
  STANDARD.
\item
  \(\rho_7\): The energy density of the matter of the fold above its
  tension in the seven-dimensional normalisation (minus the component
  \(\mathsf{t}^{(4)}{}_{(4)}\): an energy per unit volume of the six
  other directions of the fold); \(V_*\rho_7\) in the four-dimensional
  normalisation. Unit: GeV\(^7\). Status: DERIVED.
\item
  \(q_7\): The pressure along the extra times of the matter of the fold
  above its tension in the seven-dimensional normalisation (the
  component \(\mathsf{t}^{(r)}{}_{(r)}\)). Unit: GeV\(^7\). Status:
  DERIVED.
\item
  \(q_\chi\): The pressure of the particles \(\chi\) of dirac16complex
  along the extra times, integrated across the fold (their on-shell
  Lagrangian, the self-interaction energy per volume; zero for a dilute
  gas). Unit: GeV\(^4\). Status: DERIVED.
\item
  \(q_{\mathrm{P}}\): The pressure of the portals along the extra times:
  their mean on-shell Lagrangian density. Unit: GeV\(^4\). Status:
  DERIVED.
\item
  \(q_{\mathrm{SM}}\): The pressure of the Standard Model along the
  extra times at the fold: the mean on-shell Lagrangian density of the
  Standard-Model fields (the extra-time component of the energy-momentum
  tensor of every term of the fold is its Lagrangian, because the fields
  of the fold do not depend on the extra times). Unit: GeV\(^4\).
  Status: DERIVED (for radiation and free fields), ASSUMED (for the
  baryons).
\item
  \(q\): The local decay rate of a profile of dirac16complex along
  \(\sigma\) at the frequency \(\omega\),
  \(q = \sqrt{m_\Psi^2 - \omega^2\mathrm{e}^{2H\lvert\sigma\rvert}}\)
  (imaginary beyond the turning point, where the profile oscillates).
  Unit: GeV. Status: DERIVED.
\end{itemize}

\end{document}
